% !TeX encoding = UTF-8
\SetChapterInfo{subtitle={局部凸拓扑、函数空间与分布的共同基础},keywords={半范数; 拟范数; 完备化; 分布},prereq={线性代数、度量空间、Lebesgue 积分},goals={用半范数组织拓扑; 掌握连续性与完备性; 建立分布的基本运算}}
\chapter{半范数}\label{ch:yosida-01}
\ChapterSource{本章以吉田耕作第六版中译本第一章（印刷页 23--70）为主线，以英文第六版第一章（印刷页 23--67）辅助核对，保留原书全部十四节的主题与顺序。以下是统一记号后重新组织的讲述，定理编号属于本稿；标为“补充”的论证、辨析及习题用于展开原书略写的步骤。参见 \cite{yosida-zh,yosida-en}。}

半范数的作用不只是给向量赋予一个“长度”。在函数空间中，函数大小、各阶导数大小以及不同区域上的行为，往往需要分别测量。有限个或无限个这样的测量共同确定收敛，连续算子则把一种测量控制在另一种测量之内。本章从这一观点出发，依次建立局部凸空间、赋范与拟赋范空间、完备化和商空间，再把这些工具用于广义函数。

除特别说明外，标量域 $\mathbb K$ 是 $\R$ 或 $\mathbb{C}$，所有拓扑向量空间均为 Hausdorff 空间。$\Omega\subset\R^n$ 是开集，$K\Subset\Omega$ 表示 $K$ 是包含于 $\Omega$ 的紧集。多重指标 $\alpha=(\alpha_1,\ldots,\alpha_n)$ 满足 $|\alpha|=\sum\alpha_j$、$\alpha!=\prod\alpha_j!$，并记 $\partial^\alpha=\prod\partial_{x_j}^{\alpha_j}$。分布配对对测试函数是线性的，不加入复共轭；内积则约定对第一个变量线性。

\section{半范数与局部凸线性拓扑空间}\label{sec:yosida-01-01}
\noindent\textit{原书位置：中译本第 23--30 页；英文版第 23--30 页。}

\subsection{半范数及其几何}
\begin{definition}[半范数、均衡集与吸收集]
线性空间 $X$ 上的半范数是满足
\[
 p(x+y)\le p(x)+p(y),\qquad p(ax)=|a|p(x)
\]
的函数 $p:X\to[0,\infty)$。集合 $U\subset X$ 称为均衡的，若 $|a|\le1$ 蕴含 $aU\subset U$；称为吸收的，若对每个 $x\in X$，存在 $t>0$ 使 $x\in tU$。均衡凸集亦称绝对凸集。
\end{definition}
半范数可以在非零向量处为零。例如在 $\R^2$ 上，$p(x_1,x_2)=|x_1|$ 不测量第二个坐标。它的单位球是一条无限长的竖直带，而不是有界圆盘。

\begin{proposition}[基本不等式与半范数球]\label{prop:y1-seminorm-basic}
半范数满足 $p(0)=0$、$p(-x)=p(x)$ 以及
\[
 |p(x)-p(y)|\le p(x-y).
\]
若 $c>0$，则 $U_c=\{x:p(x)\le c\}$ 是吸收的绝对凸集，且
\[
 p(x)=c\inf\{t>0:x\in tU_c\}.
\]
半范数的零空间 $\ker p$ 是线性子空间。
\end{proposition}
\begin{proof}
齐次性给出前两个等式。由 $p(x)\le p(x-y)+p(y)$ 及交换 $x,y$ 后的不等式得到第三式。若 $x,y\in U_c$、$0\le s\le1$，则 $p(sx+(1-s)y)\le sc+(1-s)c=c$；均衡性由齐次性得到。选择 $t>p(x)/c$ 即得吸收性。最后 $x\in tU_c$ 等价于 $p(x)\le tc$，对 $t$ 取下确界即得公式。零空间的封闭于加法和数乘也直接来自两条公理。
\end{proof}

\begin{definition}[Minkowski 泛函]
若 $U$ 是吸收的绝对凸集，定义
\[
 p_U(x)=\inf\{t>0:x\in tU\}.
\]
这称为 $U$ 的 Minkowski 泛函或规范泛函。
\end{definition}
\begin{theorem}[Minkowski 泛函是半范数]\label{thm:y1-gauge}
$p_U$ 是有限值半范数，且
\[
 \{p_U<1\}\subset U\subset\{p_U\le1\}.
\]
若 $U$ 是拓扑向量空间中零点的绝对凸邻域，则 $p_U$ 连续。若 $U$ 还开，则 $U=\{p_U<1\}$；若 $U$ 还闭，则 $U=\{p_U\le1\}$。
\end{theorem}
\begin{proof}
吸收性保证下确界有限。对 $a\ne0$，均衡性蕴含 $aU=|a|U$，故 $p_U(ax)=|a|p_U(x)$。给定 $r>p_U(x)$、$s>p_U(y)$，均衡性保证 $x/r,y/s\in U$。凸性给出
\[
 \frac{x+y}{r+s}=\frac r{r+s}\frac xr+\frac s{r+s}\frac ys\in U.
\]
令 $r\downarrow p_U(x)$、$s\downarrow p_U(y)$ 得次可加性。若 $p_U(x)<1$，选 $r<1$ 使 $x\in rU\subset U$；另一包含关系显然。$U$ 为邻域时，$(\varepsilon/2)U\subset\{p_U<\varepsilon\}$，所以 $p_U$ 在零点连续，反三角不等式保证处处连续。

若 $U$ 开且 $x\in U$，数乘连续性给出某个 $r<1$ 使 $x/r\in U$，于是 $p_U(x)<1$。若 $U$ 闭且 $p_U(x)\le1$，则 $x/(1+1/j)\in U$，取极限得 $x\in U$。这也说明一般不能把两个包含关系不加条件地写成等式。
\end{proof}

\subsection{由半范数族生成拓扑}
\begin{theorem}[局部凸空间的半范数刻画]\label{thm:y1-lc}
设 $\mathcal P$ 是 $X$ 上一族半范数，且满足分离条件：每个 $x\ne0$ 都有 $p\in\mathcal P$ 使 $p(x)>0$。以
\[
 U(p_1,\ldots,p_r;\varepsilon_1,\ldots,\varepsilon_r)
 =\{x:p_j(x)<\varepsilon_j,\ 1\le j\le r\}
\]
为零点的邻域基，以平移得到各点邻域，则 $X$ 成为 Hausdorff 局部凸拓扑向量空间，各 $p\in\mathcal P$ 连续。反过来，每个 Hausdorff 局部凸拓扑向量空间的拓扑均由它的连续半范数生成。
\end{theorem}
\begin{proof}
上述集合的有限交仍具有同样形式，并且均为吸收的绝对凸集。把各阈值缩小一半得到 $V+V\subset U$，因而加法连续。在 $(a_0,x_0)$ 处，利用
\[
 p_j(ax-a_0x_0)\le |a|p_j(x-x_0)+|a-a_0|p_j(x_0)
\]
同时使 $a$ 接近 $a_0$、$x$ 接近 $x_0$，即可验证数乘的联合连续性。各半范数连续来自命题 \ref{prop:y1-seminorm-basic}。若 $x\ne y$，选 $p(x-y)=c>0$，则 $x+\{p<c/3\}$ 和 $y+\{p<c/3\}$ 不相交，故拓扑 Hausdorff。

反向证明中，每个零邻域包含绝对凸零邻域 $U$。定理 \ref{thm:y1-gauge} 给出连续半范数 $p_U$，且 $\{p_U<1\}\subset U$，所以这些半范数生成原拓扑。Hausdorff 性保证它们分离点，否则一个非零向量将落在所有零邻域中。
\end{proof}

图 \ref{fig:y1-seminorm} 比较单个半范数与分离点的半范数族；带状邻域的无界方向正是被忽略的方向。
\begin{figure}[htbp]
\centering
\begin{tikzpicture}[fa diagram]
\begin{scope}
  \fill[LogicColor!12] (-.8,-1.7) rectangle (.8,1.7);
  \draw[fa main,dashed] (-.8,-1.7)--(-.8,1.7) (.8,-1.7)--(.8,1.7);
  \draw[fa axis] (-2,0)--(2,0) node[right] {$x_1$};
  \draw[fa axis] (0,-1.9)--(0,2) node[above] {$x_2$};
  \draw[fa accent] (0,-1.7)--(0,1.7);
  \node[fa note] at (0,-2.5) {$p_1(x)=|x_1|<1$\\竖直方向是 $\ker p_1$};
\end{scope}
\begin{scope}[xshift=6.5cm]
  \fill[LogicColor!12] (-.8,-.8) rectangle (.8,.8);
  \draw[fa main,dashed] (-.8,-.8) rectangle (.8,.8);
  \draw[fa axis] (-2,0)--(2,0) node[right] {$x_1$};
  \draw[fa axis] (0,-1.9)--(0,2) node[above] {$x_2$};
  \node[fa note] at (0,-2.5) {$p_1(x)<1,\ p_2(x)=|x_2|<1$\\$\ker p_1\cap\ker p_2=\{0\}$};
\end{scope}
\end{tikzpicture}
\caption{一个半范数可以忽略非零方向；有限个半范数球的交形成基本零邻域。本图采用 $\R^2$ 的坐标例子，虚线边界不属于开球；一般局部凸空间不一定存在有界零邻域。}
\label{fig:y1-seminorm}
\end{figure}


把半范数族扩充为一切有限最大值及正倍数，不改变拓扑。于是许多估计可写成 $q(Tx)\le Cp(x)$；若原来的族没有经过这种扩充，则右边必须允许有限个半范数的最大值。

\subsection{光滑函数空间与测试函数空间}
函数的支集定义为其非零点集在定义域中的闭包。记
\[
 \mathcal E^k(\Omega)=C^k(\Omega),\qquad
 p_{K,m}(f)=\max_{|\alpha|\le m}\sup_{x\in K}|\partial^\alpha f(x)|,
\]
其中 $m\le k$；$k=\infty$ 时允许任意有限的 $m$。这些半范数定义局部 $C^k$ 拓扑：$f_j\to f$ 就是允许的各阶导数在每个紧集上一致收敛。取紧集穷竭 $K_j\subset\operatorname{int}K_{j+1}$，$\bigcup K_j=\Omega$，所有紧集都包含于某个 $K_j$，于是可用可数个 $p_{K_j,m}$ 定义同一拓扑。度量化将在下一节证明。

光滑紧支集函数并非只有零函数。例如
\[
 b(x)=\begin{cases}\exp\bigl(-1/(1-|x|^2)\bigr),&|x|<1,\\0,&|x|\ge1\end{cases}
\]
属于 $C_c^\infty(\R^n)$。其边界光滑性来自：任意幂次的 $t^{-1}$ 乘 $e^{-1/t}$，在 $t\downarrow0$ 时都趋于零。

令
\[
 \Test_K(\Omega)=\{\varphi\in C_c^\infty(\Omega):\supp\varphi\subset K\},\qquad
 \Test(\Omega)=C_c^\infty(\Omega).
\]
$\Test_K$ 由 $p_{K,m}$ 赋拓扑；$\Test$ 取使所有包含映射 $\Test_K\to\Test$ 连续的最强局部凸拓扑，称为局部凸归纳极限拓扑。利用穷竭可写为严格的可数归纳极限 $\Test=\bigcup_j\Test_{K_j}$。其中 $\Test_{K_j}$ 在 $\Test_{K_{j+1}}$ 中是闭子空间，诱导拓扑等于原拓扑。

\begin{proposition}[归纳极限的半范数判别]\label{prop:y1-inductive}
$\Test$ 上的半范数 $q$ 连续，当且仅当其在每个 $\Test_K$ 上的限制连续。更一般地，线性映射 $T:\Test\to Y$（$Y$ 局部凸）连续，当且仅当所有 $T|_{\Test_K}$ 连续。
\end{proposition}
\begin{proof}
一个半范数 $q$ 的各限制连续时，把 $q$ 加入生成 $\Test$ 拓扑的半范数族，仍使各包含映射连续；最强性表明 $q$ 本来就连续。必要性由复合映射的连续性得到。对一般线性映射，只须把此结论用于 $q\circ T$，其中 $q$ 遍历 $Y$ 的连续半范数。
\end{proof}

\begin{theorem}[测试函数的序列收敛]\label{thm:y1-D-sequence}
$\varphi_j\to\varphi$ 在 $\Test(\Omega)$ 中成立，当且仅当存在共同紧集 $K\Subset\Omega$ 包含所有 $\supp\varphi_j$ 及 $\supp\varphi$，且对每个 $m$，$p_{K,m}(\varphi_j-\varphi)\to0$。
\end{theorem}
\begin{proof}
充分性由包含映射连续得到。证明必要性时先令极限为零。若不存在共同紧集，可逐次选子列 $\varphi_{j_l}$ 和无域内聚点的 $x_l\in\Omega$，使 $\varphi_{j_l}(x_l)\ne0$；这是因为其支集不断逃出紧集穷竭，非零点可在逃出的支集附近选取。定义
\[
 q(\psi)=\sum_{l=1}^\infty
 \frac{|\psi(x_l)|}{|\varphi_{j_l}(x_l)|}.
\]
对任意紧支集 $\psi$ 此和只有有限项；在每个 $\Test_K$ 上，它是有限个点评价绝对值之和，所以由命题 \ref{prop:y1-inductive} 连续。然而 $q(\varphi_{j_l})\ge1$，矛盾。共同紧集存在后，全局半范数 $\psi\mapsto\sup_\Omega|\partial^\alpha\psi|$ 的各层限制连续，故在 $\Test$ 中连续，得到全部导数的一致收敛。一般极限通过相减化为零极限。
\end{proof}

图 \ref{fig:y1-test-support} 解释为什么局部光滑收敛不能替代共同紧支集条件。
\begin{figure}[htbp]
\centering
\begin{tikzpicture}[fa diagram,x=1cm,y=1cm]
  \fill[LogicColor!8] (-1.2,-.2) rectangle (1.2,1.55);
  \node at (0,1.85) {固定紧集 $K$};
  \draw[fa axis] (-2,0)--(8,0) node[right] {$x$};
  \foreach \a/\h/\c in {0/1/LogicColor,3/.65/ExColor,6/.35/DefColor}{
    \draw[thick,draw=\c] plot[domain=-1:1,samples=40] ({\a+\x},{\h*exp(1-1/(1.001-\x*\x))});
  }
  \node[fa note] at (3,-.65) {$\eta(x)$\hspace{1.2cm}$\eta(x-j)$\hspace{1cm}$a_j\eta(x-j)$};
  \draw[->,fa accent] (2.1,1.2)--(6.8,1.2) node[midway,above] {支集逃向无穷远};
\end{tikzpicture}
\caption{测试函数收敛需要共同紧支集。对非零 $\eta\in\Test(\R)$，平移在每个固定紧集上最终为零，故在 $\mathcal E$ 中趋零；即使 $a_j>0$ 趋零，$a_j\eta(\cdot-j)$ 仍不在 $\Test$ 中趋零。图中衰减高度只作示意，不给出统一导数估计。}
\label{fig:y1-test-support}
\end{figure}


\begin{pitfall}[局部收敛与测试函数收敛]
若 $\eta\ne0$ 是紧支集光滑函数，则 $\eta(x-j)$ 在 $\mathcal E(\R)$ 中趋于零，却不在 $\Test(\R)$ 中趋于零；即使乘上 $e^{-j^2}$ 也一样。共同紧支集是额外条件。$\Test$ 通常不可度量，不能把序列判别未经证明地用于任意非线性映射的连续性。
\end{pitfall}

\subsection{正则化与逼近}
\begin{proposition}[紧支集连续函数的光滑逼近]\label{prop:y1-mollify}
取非负 $\rho\in\Test(\R^n)$ 满足 $\int\rho=1$、$\supp\rho\subset\overline B(0,1)$，并令 $\rho_\varepsilon(x)=\varepsilon^{-n}\rho(x/\varepsilon)$。若 $f\in C_c(\R^n)$，则
\[
 f_\varepsilon(x)=\int f(y)\rho_\varepsilon(x-y)\,dy
\]
属于 $C_c^\infty$，支集包含于 $\supp f+\overline B(0,\varepsilon)$，且 $\|f_\varepsilon-f\|_\infty\to0$。
\end{proposition}
\begin{proof}
对光滑核在积分号下求导，得 $\partial^\alpha f_\varepsilon=f*\partial^\alpha\rho_\varepsilon$；积分区间可固定在 $\supp f$，故求导合法。支集结论来自卷积核仅在距离不超过 $\varepsilon$ 时非零。紧支集连续函数在整个 $\R^n$ 上一致连续，故
\[
 |f_\varepsilon(x)-f(x)|\le
 \sup_{|h|\le\varepsilon}|f(x-h)-f(x)|\longrightarrow0
\]
一致成立。若 $\supp f\Subset\Omega$，充分小的 $\varepsilon$ 保证 $f_\varepsilon\in\Test(\Omega)$。
\end{proof}

\begin{exercise}[补充：半范数的零空间]
证明 $X/\ker p$ 上 $\|x+\ker p\|=p(x)$ 是范数；再说明一个半范数族不分离点时，怎样先取公共零空间的商而得到 Hausdorff 空间。
\end{exercise}
\begin{solution}
若 $x-y\in\ker p$，反三角不等式给出 $p(x)=p(y)$，所以定义良好，正定性在商空间上成立，其余公理由 $p$ 继承。对半范数族令 $N=\bigcap_p\ker p$，所有 $p$ 都下降为 $X/N$ 上的半范数，且下降后的族分离商空间中的点。
\end{solution}

\section{范数和拟范数}\label{sec:yosida-01-02}
\noindent\textit{原书位置：中译本第 31--33 页；英文版第 30--32 页。}

\subsection{范数与强收敛}
\begin{definition}[赋范线性空间]
半范数 $\|\cdot\|$ 若满足 $\|x\|=0\Rightarrow x=0$，称为范数。它诱导平移不变距离 $d(x,y)=\|x-y\|$。在此距离中的收敛称为范数收敛或强收敛，记作 $x_j\to x$，原书也用 $\mathrm{s}\!\operatorname{-lim}x_j=x$。
\end{definition}
范数满足反三角不等式；若 $x_j\to x$、$y_j\to y$、$a_j\to a$，则
\[
 \|x_j\|\to\|x\|,\qquad x_j+y_j\to x+y,\qquad a_jx_j\to ax.
\]
前两条由三角不等式得到，最后一条由
\[
 \|a_jx_j-ax\|\le |a_j|\|x_j-x\|+|a_j-a|\|x\|
\]
及收敛标量序列有界得到。这是在最简单情形下的拓扑与线性运算相容性。

\subsection{吉田意义的拟范数}
\begin{definition}[拟范数：本章的历史用法]\label{def:y1-Fnorm}
函数 $q:X\to[0,\infty)$ 若满足
\[
 q(x)=0\Longleftrightarrow x=0,\quad q(-x)=q(x),\quad
 q(x+y)\le q(x)+q(y),
\]
并满足两条数乘连续性条件
\[
 a_j\to0\ \Longrightarrow\ q(a_jx)\to0\quad(x\text{ 固定}),\qquad
 q(x_j)\to0\ \Longrightarrow\ q(ax_j)\to0\quad(a\text{ 固定}),
\]
称为吉田意义的拟范数。它定义距离 $d(x,y)=q(x-y)$。
\end{definition}
\begin{remark}[补充：与现代术语的区别]
原书的 quasi-norm 在这里不要求 $q(ax)=|a|q(x)$，但要求真正的三角不等式。现代文献常以 quasi-norm 指满足绝对齐次性和放宽三角不等式 $\|x+y\|\le C(\|x\|+\|y\|)$ 的函数。为避免混用，本章在谈完备度量线性空间时使用定义 \ref{def:y1-Fnorm}，也可称之为一种 F-范数。
\end{remark}

\begin{lemma}[Kakutani 的一致数乘结论]\label{lem:y1-uniform-scalars}
若 $q(x_j)\to0$，则对每个 $M<\infty$，
\[
 \sup_{|a|\le M}q(ax_j)\longrightarrow0.
\]
因此定义 \ref{def:y1-Fnorm} 中数乘实际上联合连续。
\end{lemma}
\begin{proof}
这是原书借助 Egorov 定理给出的关键步骤，现将其展开。令 $r_j(a)=q(ax_j)$。由反三角不等式与第一条数乘条件，各 $r_j$ 在标量域 $\mathbb K$ 上连续；第二条条件给出对每个 $a$，$r_j(a)\to0$。在标量域的单位球上应用 Egorov 定理，可取有限正 Lebesgue 测度的可测集 $A$，使 $r_j\to0$ 在 $A$ 上一致。复标量时把 $\mathbb{C}$ 看作 $\R^2$。

可测集示性函数的 $L^1$ 平移连续性给出 $|A\mathbin\triangle(A+h)|\to0$。这里这一事实可由有限测度集在测度上用有限个矩形逼近、再利用矩形的平移连续性证明。因此存在 $\delta>0$，使 $|h|<\delta$ 时 $A\cap(A+h)\ne\varnothing$，即 $h=a-b$，其中 $a,b\in A$。于是
\[
 r_j(h)\le r_j(a)+r_j(-b)=r_j(a)+r_j(b)
 \le2\sup_{c\in A}r_j(c).
\]
这给出 $r_j\to0$ 在 $|h|<\delta$ 上一致。给定 $M$，选整数 $N>M/\delta$，对 $|a|\le M$ 写 $a=N(a/N)$，次可加性给出 $r_j(a)\le Nr_j(a/N)$，即得结论。

若 $a_j\to a$、$x_j\to x$，则
\[
 q(a_jx_j-ax)\le q(a_j(x_j-x))+q((a_j-a)x)\to0.
\]
第一项使用刚证的一致结论，第二项使用固定向量的数乘条件。加法与 $q$ 自身的连续性则来自三角和反三角不等式。
\end{proof}

\subsection{可数半范数族的度量化}
\begin{theorem}[可数半范数族与平移不变距离]\label{thm:y1-countable-metric}
设分离点的半范数族为 $p_1,p_2,\ldots$。则
\[
 q(x)=\sum_{j=1}^\infty2^{-j}\frac{p_j(x)}{1+p_j(x)},\qquad d(x,y)=q(x-y)
\]
定义吉田意义的拟范数与距离，生成同一局部凸拓扑。并且 $d$-Cauchy 等价于对每个 $j$ 在 $p_j$ 下 Cauchy。
\end{theorem}
\begin{proof}
函数 $h(t)=t/(1+t)$ 增加、$h(0)=0$，且
\[
 h(s+t)=\frac{s}{1+s+t}+\frac{t}{1+s+t}\le h(s)+h(t).
\]
因此各求和项满足三角不等式，级数由 $\sum2^{-j}$ 控制。分离性给出正定性。若 $q(x_l)\to0$，则每个固定项趋于零，从而 $p_j(x_l)\to0$。反向结论先固定足够小的级数尾，再控制有限个首项得到。此“有限首项加一致小尾项”论证也证明两个邻域系统一致，以及两种 Cauchy 条件等价。半范数的齐次性与同一论证验证两条数乘条件。
\end{proof}
于是 $\mathcal E^k(\Omega)$、$\Test_K(\Omega)$ 可度量。不同的穷竭或等价半范数族给出不同公式的距离，但拓扑与 Cauchy 概念相同。

\begin{proposition}[补充：可赋范性的判别]\label{prop:y1-normable}
Hausdorff 局部凸空间可由一个范数生成拓扑，当且仅当它有一个有界的零邻域。
\end{proposition}
\begin{proof}
范数单位球显然是有界零邻域。反向取有界零邻域中的绝对凸零邻域 $U$。其规范泛函 $p_U$ 连续。若 $p_U(x)=0$，则 $x\in tU$ 对每个 $t>0$ 成立；有界性使 $tU$ 最终包含于任意给定零邻域，故 $x=0$。对任意零邻域 $V$，有界性给出 $U\subset cV$，于是 $\{p_U<1/c\}\subset V$（可先取均衡 $V$）。所以 $p_U$ 的范数拓扑等于原拓扑。
\end{proof}
有界集的正式定义见第 \ref{sec:yosida-01-07} 节。此判别说明可度量性远弱于可赋范性：多个局部测量不能总由一种全局测量替代。

\section{几例赋范线性空间}\label{sec:yosida-01-03}
\noindent\textit{原书位置：中译本第 33--39 页；英文版第 32--38 页。}

\subsection{有界连续函数与可积函数}
拓扑空间 $S$ 上的有界连续函数组成 $C_b(S)$，赋予范数 $\|f\|_\infty=\sup_S|f|$。范数收敛就是全域一致收敛。原书此处的 $C(S)$ 限于有界连续函数；若 $S$ 紧，所有连续函数都自动有界，记号 $C(S)$ 无歧义。

对测度空间 $(S,\Sigma,\mu)$ 与 $1\le p<\infty$，令
\[
 L^p(S,\mu)=\{f\text{ 可测}:\int_S|f|^p\,d\mu<\infty\}/\text{几乎处处相等},
 \qquad\|f\|_p=\left(\int_S|f|^p\,d\mu\right)^{1/p}.
\]
商掉几乎处处相等是正定性的必要条件。$L^\infty$ 由本质有界函数的等价类组成，范数为
\[
 \|f\|_\infty=\operatorname*{ess\,sup}|f|
 =\inf\{M\ge0:|f|\le M\text{ 几乎处处}\}.
\]
本质上确界与逐点上确界不同：零测集上的任意修改都不改变前者。

\begin{lemma}[Young 不等式]
若 $1<p<\infty$、$1/p+1/p'=1$，则对 $a,b\ge0$，
\[
 ab\le\frac{a^p}{p}+\frac{b^{p'}}{p'}.
\]
等号成立当且仅当 $a^p=b^{p'}$。
\end{lemma}
\begin{proof}
固定 $b$，函数 $a\mapsto a^p/p-ab+b^{p'}/p'$ 的导数为 $a^{p-1}-b$，唯一极小值点为 $a=b^{1/(p-1)}$，极小值为零。$b=0$ 的情形直接验证。
\end{proof}
\begin{theorem}[Hölder 与 Minkowski 不等式]\label{thm:y1-holder}
对共轭指数 $p,p'$（包含 $p=1,p'=\infty$），
\[
 \int|fg|\,d\mu\le\|f\|_p\|g\|_{p'}.
\]
对 $1\le p\le\infty$，$\|f+g\|_p\le\|f\|_p+\|g\|_p$，因而 $L^p$ 是赋范线性空间。
\end{theorem}
\begin{proof}
端点 Hölder 由 $|g|\le\|g\|_\infty$ 几乎处处得到。内部指数时，若两范数不为零，将 Young 不等式用于 $|f|/\|f\|_p$ 和 $|g|/\|g\|_{p'}$ 后积分；任一范数为零则直接成立。

先由 $|f+g|^p\le2^{p-1}(|f|^p+|g|^p)$ 确保 $f+g\in L^p$。当 $1<p<\infty$，利用 Hölder 和 $(p-1)p'=p$，
\[
 \|f+g\|_p^p
 \le\int |f+g|^{p-1}(|f|+|g|)\,d\mu
 \le\|f+g\|_p^{p-1}(\|f\|_p+\|g\|_p).
\]
若左侧非零，约去相应因子即得 Minkowski；否则结论显然。$p=1,\infty$ 直接由逐点三角不等式得到。
\end{proof}
当 $1<p<\infty$ 且 $f,g$ 非零时，Minkowski 等号成立恰在 $f=cg$ 几乎处处、$c>0$ 的情形。实值时是同号正比例，复值时还要求相同相位；证明由 Hölder 的模值比例条件与逐点三角等号条件合并得到。$p=1$ 只须几乎处处无相消，不要求固定比例；$p=\infty$ 的等号条件也不同。

\begin{proposition}[有限测度空间上的极限]
若 $\mu(S)<\infty$ 且 $f\in L^\infty$，则 $\lim_{p\to\infty}\|f\|_p=\|f\|_\infty$。
\end{proposition}
\begin{proof}
总测度为零或 $\|f\|_\infty=0$ 时直接成立。记 $M=\|f\|_\infty>0$，上界为 $\|f\|_p\le\mu(S)^{1/p}M$。对 $0<\varepsilon<M$，集合 $A_\varepsilon=\{|f|>M-\varepsilon\}$ 有正测度，所以
\[
 \|f\|_p\ge\mu(A_\varepsilon)^{1/p}(M-\varepsilon).
\]
先令 $p\to\infty$，再令 $\varepsilon\downarrow0$。
\end{proof}

\subsection{数列空间}
在 $\N$ 上取计数测度，即得
\[
 \ell^p=\{x=(x_j):\sum|x_j|^p<\infty\},\qquad
 \ell^\infty=\{x:\sup_j|x_j|<\infty\}.
\]
另外 $c_0=\{x:x_j\to0\}$、$c=\{x:x_j\text{ 有有限极限}\}$ 均取上确界范数。它们是 $\ell^\infty$ 的闭子空间：一致极限保留趋零或存在极限的条件。有限支集数列在 $c_0$ 与 $\ell^p$（$p<\infty$）中稠密，却不在 $\ell^\infty$ 中稠密，常数数列 $(1,1,\ldots)$ 与任何有限支集数列距离都是 $1$。

\subsection{带符号测度、Jordan 分解与 Hahn 分解}
以下先取 $\sigma$-代数 $\Sigma$。原书用 $\sigma$-环表述的情形可对可测集合分别限制处理；为使全空间上的公式直接有意义，此处假定 $S\in\Sigma$。一个有限值带符号测度 $\nu:\Sigma\to\R$ 是可数可加的集合函数。有限值指每个可测集的值都是有限实数，随后将证明其全变差有限。

\begin{lemma}[有限值带符号测度的值域有界]\label{lem:y1-measure-bounded}
$\sup_{A\in\Sigma}|\nu(A)|<\infty$。
\end{lemma}
\begin{proof}
记 $V(A)=\sup_{B\subset A}|\nu(B)|$。若 $V(S)=\infty$，归纳构造递减 $E_j$，使 $V(E_j)=\infty$ 且 $|\nu(E_j)|\ge j$。给定 $E_j$，选 $B\subset E_j$，$|\nu(B)|>|\nu(E_j)|+j+1$。$B$ 与 $E_j\setminus B$ 的测度绝对值都大于 $j+1$，并且至少一者有无穷大的 $V$；选该者为 $E_{j+1}$。可是可数可加性应用于不交分解
\[
 E_1=\left(\bigcap_jE_j\right)\mathbin{\dot\cup}
 \bigcup_j(E_j\setminus E_{j+1})
\]
给出 $\nu(E_j)\to\nu(\bigcap_jE_j)$，与绝对值无界矛盾。
\end{proof}

\begin{theorem}[Jordan 与 Hahn 分解]\label{thm:y1-jordan}
定义
\[
 \nu^+(A)=\sup_{B\subset A}\nu(B),\qquad
 \nu^-(A)=-\inf_{B\subset A}\nu(B).
\]
则 $\nu^+,\nu^-$ 是有限正测度，$\nu=\nu^+-\nu^-$。存在可测分解 $S=P\mathbin{\dot\cup}N$，使 $P$ 的每个可测子集测度非负，$N$ 的每个可测子集测度非正；而且 $\nu^-(P)=0$、$\nu^+(N)=0$。这种分解称为 Hahn 分解。
\end{theorem}
\begin{proof}
上一引理保证上下确界有限。对不交 $A_j$，任意 $B\subset\bigcup A_j$ 有
\[
 \nu(B)=\sum_j\nu(B\cap A_j)\le\sum_j\nu^+(A_j),
\]
故 $\nu^+(\bigcup A_j)\le\sum\nu^+(A_j)$。反向对前 $m$ 个 $A_j$ 各取逼近上确界的子集，将它们合并，误差取小于任意 $\varepsilon$；再令 $\varepsilon\downarrow0$、$m\to\infty$。这证明可数可加性。对 $-\nu$ 同理。由 $\nu(B)=\nu(A)-\nu(A\setminus B)$ 对 $B\subset A$ 取上确界，得 $\nu^+(A)=\nu(A)+\nu^-(A)$。

选 $B_j$ 使 $\nu(B_j)>\nu^+(S)-2^{-j}$。由分解公式得 $\nu^-(B_j)<2^{-j}$、$\nu^+(S\setminus B_j)<2^{-j}$。令 $P=\liminf B_j$、$N=S\setminus P$。对每个 $m$，
\[
 N\subset\bigcup_{j\ge m}(S\setminus B_j),\qquad
 \nu^+(N)\le\sum_{j\ge m}2^{-j}\to0.
\]
对 $P_m=\bigcap_{j\ge m}B_j$，有 $\nu^-(P_m)\le\inf_{j\ge m}2^{-j}=0$；$P=\bigcup_mP_m$，故 $\nu^-(P)=0$。这正是所需的符号性质。
\end{proof}
原书以负值的“负变差”记下确界；此处 $\nu^-$ 取非负值，因此 Jordan 公式使用减号。这个记号改变不改变数学内容。

\begin{proposition}[全变差范数]\label{prop:y1-variation}
令 $|\nu|=\nu^++\nu^-$，则
\[
 \|\nu\|_{\mathrm{TV}}=|\nu|(S)
 =\sup_{\substack{h\text{ 可测}\\|h|\le1}}
 \left|\int h\,d\nu\right|
 =\sup_{\{A_j\}}\sum_j|\nu(A_j)|,
\]
最后的上确界取遍有限可测分割。有限带符号测度组成赋范实线性空间。有限复测度也组成赋范复线性空间，其范数由同一分割公式给出，积分上确界中允许复值 $h$。
\end{proposition}
\begin{proof}
实情形中 $|\int h\,d\nu|\le\int|h|\,d|\nu|\le|\nu|(S)$；取 $h=1_P-1_N$ 达到上界。每个分割的和不超过全变差，而分割 $P,N$ 达到上界。复情形中 $\nu=\nu_1+i\nu_2$，由前述实测度结论知分割和不超过 $|\nu_1|(S)+|\nu_2|(S)<\infty$。对一个有限分割，取简单函数在 $A_j$ 上为 $\overline{\nu(A_j)}/|\nu(A_j)|$（零值时取零），积分等于该分割和。反向对简单 $h$ 逐项估计，再用一致简单函数逼近一般有界可测 $h$，得积分公式。两种情形下范数的三角性及齐次性来自积分公式；若范数为零，取 $h=1_A$ 得 $\nu(A)=0$ 对所有 $A$ 成立。
\end{proof}

\begin{exercise}[补充：全变差不是总质量的绝对值]
在两点空间取 $\nu=\delta_a-\delta_b$。计算 $\nu(S)$ 和 $\|\nu\|_{\mathrm{TV}}$，并用积分上确界解释两者为什么不同。
\end{exercise}
\begin{solution}
总质量为零而全变差为二。选择 $h(a)=1,h(b)=-1$，则 $\int h\,d\nu=2$；全变差可以分别选择不同位置的符号，避免正负质量抵消。
\end{solution}

\section{几例拟赋范线性空间}\label{sec:yosida-01-04}
\noindent\textit{原书位置：中译本第 39--40 页；英文版第 38--39 页。}

\subsection{局部光滑拓扑}
按定理 \ref{thm:y1-countable-metric}，$\mathcal E^k(\Omega)$ 和固定紧集上的 $\Test_K$ 都可取距离
\[
 d(f,g)=\sum_{j=1}^\infty2^{-j}\frac{p_j(f-g)}{1+p_j(f-g)},
\]
其中 $p_j$ 枚举相应的半范数。此拟范数可有统一上界 $1$，但这不意味着全空间是拓扑有界集：拓扑有界性涉及标量缩放与所有邻域，而不是某个距离数值的上界。

\subsection{有限测度空间上的测度收敛}
\begin{definition}[可测函数空间]
设 $\mu(S)<\infty$。用 $L^0(S,\mu)$（原书记 $M(S,\Sigma,\mu)$）表示几乎处处有限的可测函数按几乎处处相等组成的空间。定义
\[
 q(f)=\int_S\frac{|f|}{1+|f|}\,d\mu,\qquad d(f,g)=q(f-g).
\]
\end{definition}
\begin{proposition}[积分距离与测度收敛]\label{prop:y1-measure-convergence}
上述 $q$ 是吉田意义的拟范数，并且 $d(f_j,f)\to0$ 当且仅当对每个 $\varepsilon>0$，
\[
 \mu\{|f_j-f|\ge\varepsilon\}\longrightarrow0.
\]
\end{proposition}
\begin{proof}
正定性在等价类上成立。三角不等式使用 $t/(1+t)$ 的单调性与次可加性。固定 $f$ 时，$q(af)\to0$ 随 $a\to0$ 由控制收敛得到，控制函数是可积的常数 $1$。

令 $A_\varepsilon=\{|u|\ge\varepsilon\}$，则
\[
 \frac{\varepsilon}{1+\varepsilon}\mu(A_\varepsilon)
 \le q(u)\le\mu(A_\varepsilon)
 +\frac{\varepsilon}{1+\varepsilon}\mu(S).
\]
第一式证明距离收敛推出测度收敛；第二式先令序列指标趋于无穷，再令 $\varepsilon\downarrow0$，证明反向。固定 $a\ne0$ 时，$\{|au|\ge\varepsilon\}=\{|u|\ge\varepsilon/|a|\}$，故第二条数乘条件成立。联合连续性也可引用引理 \ref{lem:y1-uniform-scalars}。
\end{proof}
原书还给出同一拓扑的另一种拟范数。令
\[
 r(f)=\inf_{\varepsilon>0}\bigl(\varepsilon+\mu\{|f|\ge\varepsilon\}\bigr),
 \qquad \widetilde q(f)=\arctan r(f).
\]
由 $\{|f+g|\ge a+b\}\subset\{|f|\ge a\}\cup\{|g|\ge b\}$ 得 $r(f+g)\le r(f)+r(g)$；$\arctan$ 在非负轴上增加且次可加，故 $\widetilde q$ 满足三角不等式。$r(f_j)\to0$ 与测度收敛等价：一向取固定阈值，另一向选择接近下确界的小阈值，再比较较大阈值的异常集。于是它也满足数乘条件，并定义相同拓扑。

\begin{example}[测度收敛与几乎处处收敛]
在 $[0,1)$ 上依次枚举第 $m$ 层的二进区间
\[
 I_{m,k}=[k2^{-m},(k+1)2^{-m}),\quad 0\le k<2^m,
\]
取相应示性函数作为一个序列。其积分距离趋于零，却在每个点上都无极限，因为每层恰有一次取值为 $1$，其余大量次取值为 $0$。相反，有限测度空间上的几乎处处收敛必推出测度收敛，由控制收敛应用于异常集示性函数即可。
\end{example}
测度收敛序列可以抽出几乎处处收敛的子列：选 $j_k$ 使 $\mu\{|f_{j_k}-f|>2^{-k}\}<2^{-k}$，可数次可加性给出异常事件的上极限为零测集。这一选子列技术也用于后面的完备性证明。

\begin{remark}[补充：非局部凸性需要测度条件]
在 $[0,1]$ 的 Lebesgue 测度上，$L^0$ 的测度拓扑并非局部凸。事实上其连续线性泛函只能为零：连续性给出某个 $\delta>0$，使 $q(f)<\delta$ 时 $|L(f)|<1$。任何支集测度小于 $\delta$ 的 $f$，其任意标量倍 $tf$ 都有 $q(tf)\le\mu(\supp f)<\delta$，故 $L(f)=0$。把 $[0,1]$ 分成有限个长度小于 $\delta$ 的区间，将任意 $f$ 分解于这些区间上，即得 $L(f)=0$。若此非零 Hausdorff 空间局部凸，Hahn--Banach 定理将给出分离点的连续线性泛函，矛盾。有限原子空间则可局部凸，所以不能对任意测度空间一概作此断言。
\end{remark}

\section{准 Hilbert 空间}\label{sec:yosida-01-05}
\noindent\textit{原书位置：中译本第 40--44 页；英文版第 39--42 页。}

\subsection{平行四边形恒等式与极化}
\begin{definition}[准 Hilbert 空间]
赋范空间若满足
\[
 \|x+y\|^2+\|x-y\|^2=2\|x\|^2+2\|y\|^2,
\]
称为准 Hilbert 空间，亦称预 Hilbert 空间。这里尚不要求完备。
\end{definition}
\begin{theorem}[Jordan--von Neumann 极化定理]\label{thm:y1-polarization}
实赋范空间满足平行四边形恒等式，当且仅当其范数由唯一实内积生成，此内积为
\[
 B(x,y)=\frac14\bigl(\|x+y\|^2-\|x-y\|^2\bigr).
\]
复情形中，以第一变量线性为约定，唯一内积为
\[
 (x,y)=\frac14\left(\|x+y\|^2-\|x-y\|^2
 +i\|x+iy\|^2-i\|x-iy\|^2\right).
\]
\end{theorem}
\begin{proof}
实情形中对称性及 $B(x,x)=\|x\|^2$ 直接成立。平行四边形恒等式分别用于 $x+z,y$ 和 $x-z,y$，相减得到
\[
 B(x+y,z)+B(x-y,z)=2B(x,z).
\]
取 $y=x$ 得 $B(2x,z)=2B(x,z)$。再令 $x=(u+v)/2$、$y=(u-v)/2$，得到 $B(u+v,z)=B(u,z)+B(v,z)$。于是整数及有理数齐次性成立；$B$ 的定义与范数连续性使有理数逼近推广为实齐次性。这完整证明了双线性。

复空间先视为实空间，得到 $B$。由 $\|ix\|=\|x\|$ 可知 $B(ix,iy)=B(x,y)$，所以 $B(ix,y)=-B(x,iy)$，且 $B(x,ix)=0$。定义 $(x,y)=B(x,y)+iB(x,iy)$，便有 $(ix,y)=i(x,y)$、$(y,x)=\overline{(x,y)}$ 及 $(x,x)=\|x\|^2$。这证明复内积性质。内积展开 $\|x\pm y\|^2$ 和 $\|x\pm iy\|^2$ 反过来给出极化公式，故唯一性与必要性都成立。
\end{proof}

\begin{theorem}[内积生成范数与 Cauchy--Schwarz 不等式]\label{thm:y1-schwarz}
若 $(x,y)$ 是正定的实或复内积，则 $\|x\|=(x,x)^{1/2}$ 是范数，满足平行四边形恒等式，并且
\[
 |(x,y)|\le\|x\|\|y\|.
\]
等号成立当且仅当 $x,y$ 线性相关。
\end{theorem}
\begin{proof}
$y=0$ 时显然。否则令 $c=(x,y)/\|y\|^2$，利用第一变量线性展开
\[
 0\le\|x-cy\|^2=\|x\|^2-\frac{|(x,y)|^2}{\|y\|^2}.
\]
等号恰在 $x=cy$ 时成立。因而
\[
 \|x+y\|^2=\|x\|^2+2\operatorname{Re}(x,y)+\|y\|^2
 \le(\|x\|+\|y\|)^2.
\]
这是三角不等式；其余范数公理及平行四边形恒等式直接由内积展开得到。
\end{proof}
\begin{remark}[补充：完备性与共轭位置]
仅有内积公理不能推出 Hilbert 空间，还须证明完备性；原书中译本这一处的表述应按英文版的 pre-Hilbert 理解。极化公式中虚部的符号取决于哪一个变量线性。本稿内积第一变量线性，而分布对测试函数线性，二者不应混用。
\end{remark}

\subsection{基本例子与 Sobolev 型内积}
$L^2(S,\mu)$ 和 $\ell^2$ 分别取
\[
 (f,g)=\int f\overline g\,d\mu,\qquad
 (x,y)=\sum_jx_j\overline{y_j}.
\]
积分与级数的绝对收敛由 Hölder 保证。对 $k\in\N_0$，令 $\widehat H^k(\Omega)$ 为满足以下有限范数的 $C^k$ 函数：
\[
 \|f\|_{k,2}^2=\sum_{|\alpha|\le k}\int_\Omega|\partial^\alpha f|^2\,dx,
 \qquad (f,g)_{k,2}=\sum_{|\alpha|\le k}\int_\Omega
 \partial^\alpha f\,\overline{\partial^\alpha g}\,dx.
\]
它是准 Hilbert 空间，$C_c^\infty(\Omega)$ 也可用此内积赋予准 Hilbert 结构。它们一般尚不完备；后面将用弱导数解释其完备化。原书也考虑 $C_c^k$，其完备化相同：先用正则化在全部有限阶导数范数中逼近，即可将 $C_c^k$ 换成 $C_c^\infty$。

\subsection{Bergman 空间与 Hardy--Lebesgue 类}
平面开集 $G$ 上的 Bergman 空间为
\[
 A^2(G)=\{f\text{ 在 }G\text{ 全纯}:\int_G|f|^2\,dA<\infty\},
 \qquad(f,g)=\int_Gf\overline g\,dA.
\]
它继承 $L^2$ 内积。第 \ref{sec:yosida-01-09} 节将证明全纯性在 $L^2$ 极限中保留，从而得到完备性。

单位圆盘 $\mathbb D$ 上的 Hardy--Lebesgue 类（现代记作 $H^2(\mathbb D)$）由满足
\[
 \|f\|_{H^2}^2=\sup_{0<r<1}\frac1{2\pi}
 \int_0^{2\pi}|f(re^{i\theta})|^2\,d\theta<\infty
\]
的全纯函数组成。
\begin{theorem}[Hardy 类与系数空间的等距对应]\label{thm:y1-hardy}
若 $f(z)=\sum_{j=0}^\infty c_jz^j$，则 $f\in H^2$ 当且仅当 $(c_j)\in\ell^2$，且 $\|f\|_{H^2}^2=\sum|c_j|^2$。系数映射是线性等距双射，因此内积为 $(f,g)=\sum c_j\overline{d_j}$。
\end{theorem}
\begin{proof}
在每个固定 $r<1$ 上，Taylor 级数一致绝对收敛，可逐项相乘积分。圆周指数函数的正交性给出
\[
 \frac1{2\pi}\int_0^{2\pi}|f(re^{i\theta})|^2\,d\theta
 =\sum_{j=0}^\infty|c_j|^2r^{2j}.
\]
令 $r\uparrow1$，单调收敛即得范数公式。反过来，给定 $\ell^2$ 系数，Cauchy--Schwarz 给出对 $|z|\le r<1$，
\[
 \sum_{j\ge N}|c_jz^j|
 \le\left(\sum_{j\ge N}|c_j|^2\right)^{1/2}
 \frac{r^N}{\sqrt{1-r^2}}.
\]
因此级数在紧子圆盘上一致收敛，定义全纯函数，前一公式保证其属于 $H^2$。这也证明满射性。
\end{proof}

\begin{exercise}[补充：哪些序列范数来自内积]
在 $\mathbb K^2$ 上取 $\ell^p$ 范数。用 $(1,0)$ 和 $(0,1)$ 验证：$1\le p\le\infty$ 中只有 $p=2$ 满足平行四边形恒等式。
\end{exercise}
\begin{solution}
有限 $p$ 时，左边为 $2\cdot2^{2/p}$，右边为 $4$，故 $p=2$。$p=\infty$ 时左边为 $2$，也不等于 $4$。$p=2$ 的范数由标准内积生成，故对所有向量成立。
\end{solution}

\section{线性算子的连续性}\label{sec:yosida-01-06}
\noindent\textit{原书位置：中译本第 44--46 页；英文版第 42--44 页。}

线性算子的定义域 $D(T)$ 是线性子空间，取原空间的诱导拓扑。由 $T(x+h)-Tx=Th$，在零点连续等价于在定义域每一点连续。

\begin{theorem}[连续线性算子的半范数估计]\label{thm:y1-operator-seminorm}
设 $X,Y$ 分别由半范数族 $\mathcal P,\mathcal Q$ 定义拓扑。线性映射 $T:D(T)\subset X\to Y$ 连续，当且仅当对每个 $q\in\mathcal Q$，存在 $p_1,\ldots,p_r\in\mathcal P$ 与 $C>0$，使
\[
 q(Tx)\le C\max_{1\le j\le r}p_j(x)\quad(x\in D(T)).
\]
若 $\mathcal P$ 对有限最大值封闭，则可简写为 $q(Tx)\le Cp(x)$。
\end{theorem}
\begin{proof}
连续性给出有限半范数邻域，其与定义域的交映入 $\{q<1\}$。令 $p(x)=\max_jp_j(x)/\varepsilon_j$。若 $p(x)>0$，对每个 $t>p(x)$，$x/t$ 落在该邻域中，于是 $q(Tx)<t$，令 $t\downarrow p(x)$。若 $p(x)=0$，所有倍数 $ax$ 都在该邻域中，$|a|q(Tx)<1$ 对任意 $a$ 成立，故 $q(Tx)=0$。反向估计把每个有限的目标邻域条件转成有限的源邻域条件，因此在零点连续。
\end{proof}
\begin{corollary}[线性泛函、范数估计与逆算子]
局部凸空间上的线性泛函连续，当且仅当其绝对值受有限个连续半范数的最大值控制。赋范空间之间的线性算子连续，当且仅当 $\|Tx\|\le C\|x\|$。而逆映射 $T^{-1}:R(T)\to D(T)$ 存在且连续，当且仅当存在 $c>0$，使
\[
 \|Tx\|\ge c\|x\|\quad(x\in D(T)).
\]
\end{corollary}
\begin{proof}
前两项分别以标量绝对值及单一范数作为目标半范数。后一项的下界首先使 $T$ 单射；把它改写成 $\|T^{-1}y\|\le c^{-1}\|y\|$ 即得连续性，反向同理。这里逆算子只定义于值域，尚未断言值域是整个目标空间。
\end{proof}

\begin{definition}[算子范数]
全定义的有界线性算子 $T:X\to Y$ 的算子范数为
\[
 \|T\|=\sup_{\|x\|\le1}\|Tx\|
 =\inf\{C\ge0:\|Tx\|\le C\|x\|\ \forall x\}.
\]
当 $X\ne\{0\}$ 时，还等于 $\sup_{\|x\|=1}\|Tx\|$。此处“有界”指把单位球映成有界集，等价于连续；不指在整个 $X$ 上函数值有统一上界。
\end{definition}
若 $S,T$ 定义域不同，则 $D(S+T)=D(S)\cap D(T)$，$D(ST)=\{x\in D(T):Tx\in D(S)\}$。对全定义有界算子，
\[
 \|S+T\|\le\|S\|+\|T\|,\quad
 \|aT\|=|a|\|T\|,\quad \|ST\|\le\|S\|\|T\|.
\]
这些公式逐点估计后在单位球取上确界即可；例如 $\|STx\|\le\|S\|\|Tx\|\le\|S\|\|T\|\|x\|$。于是 $\|T^m\|\le\|T\|^m$，$T^0=I$。

\begin{example}[算子的非交换性与定义域]
在 $L^2(\R)$ 上，乘法算子 $Mu(t)=tu(t)$ 的定义域为 $\{u:tu\in L^2\}$，微分算子 $Du=iu'$ 可取定义域 $H^1(\R)$。在共同的不变子空间 $\Test(\R)$ 上直接算得
\[
 (DM-MD)u=i(tu)'-tiu'=iu.
\]
这两个算子在 $L^2$ 范数下都不是全定义有界算子。给出共同定义域后，交换关系才有明确含义。
\end{example}

\begin{proposition}[补充：有界下界与闭值域]
若 $X$ 是 Banach 空间，$T:X\to Y$ 有界且 $\|Tx\|\ge c\|x\|$，则 $R(T)$ 在赋范空间 $Y$ 中闭。
\end{proposition}
\begin{proof}
若 $Tx_j\to y$，下界使 $x_j$ 是 Cauchy 序列，故趋于 $x\in X$。连续性给出 $y=Tx$。源空间的完备性在此不可省去，例如未完备空间到其完备化的等距包含有稠密但不闭的值域。
\end{proof}

\section{有界集和囿空间}\label{sec:yosida-01-07}
\noindent\textit{原书位置：中译本第 46--48 页；英文版第 44--46 页。}

\begin{definition}[拓扑有界集与囿空间]
拓扑向量空间中的 $B$ 称为有界集，若每个零邻域 $U$ 都吸收 $B$：存在 $c>0$ 使 $B\subset cU$（可先把 $U$ 换成均衡零邻域）。局部凸空间称为囿空间或 bornological 空间，若每个吸收全部有界集的绝对凸集都是零邻域。
\end{definition}
\begin{proposition}[有界集的半范数判别]\label{prop:y1-bounded}
局部凸空间中，$B$ 有界当且仅当对每个连续半范数 $p$，$\sup_{x\in B}p(x)<\infty$。连续线性映射把有界集映为有界集。紧集以及收敛序列连同其极限均有界。
\end{proposition}
\begin{proof}
有界性应用于 $\{p<1\}$ 得一向；反向对一个含有限个半范数条件的零邻域，选择足够大的共同缩放常数即可。若 $T$ 连续，给定目标零邻域 $V$，取源零邻域 $U$ 使 $T(U)\subset V$，则 $B\subset cU$ 推出 $T(B)\subset cV$。紧集的有界性由零邻域的吸收性和有限子覆盖证明；收敛序列加极限是紧集。这最后一条对一般拓扑向量空间也成立。
\end{proof}

\begin{theorem}[囿空间的等价刻画]\label{thm:y1-bornological}
局部凸空间 $X$ 是囿空间，当且仅当每个在全部有界集上有界的半范数都连续。因而从囿空间到局部凸空间的线性映射连续，当且仅当它把有界集映为有界集。
\end{theorem}
\begin{proof}
设 $X$ 囿，半范数 $p$ 在各有界集上有界。则 $\{p\le1\}$ 吸收所有有界集，故是零邻域；由半范数的缩放得 $p$ 连续。反向设每个这种半范数连续。若绝对凸集 $U$ 吸收全部有界集，它尤其吸收每个单点，因此 $p_U$ 有定义；吸收有界集保证 $p_U$ 在有界集上有界，于是连续，$\{p_U<1\}\subset U$ 证明 $U$ 是邻域。

对线性映射 $T:X\to Y$，若它保有界集，则每个连续目标半范数 $q$ 的复合 $q\circ T$ 在源有界集上有界，故连续；半范数刻画推出 $T$ 连续。反向已由上一命题证明。
\end{proof}
赋范空间都是囿空间：若半范数 $p$ 在范数单位球上有界于 $C$，则 $p(x)\le C\|x\|$，所以连续。不要求赋范空间完备。

\begin{theorem}[拟赋范空间上保有界集的线性映射]
设 $X,Y$ 均为定义 \ref{def:y1-Fnorm} 意义下的拟赋范空间。线性映射 $T:X\to Y$ 连续，当且仅当它保拓扑有界集。
\end{theorem}
\begin{proof}
只证明充分性。若 $x_j\to0$，可选正整数 $n_j\to\infty$，使 $n_jx_j\to0$：对每个整数 $k$，固定数乘的连续性给出 $kx_j\to0$，选择递增的 $N_k$ 使 $j\ge N_k$ 时 $q_X(kx_j)<1/k$，再令 $n_j=k$ 对 $N_k\le j<N_{k+1}$。于是 $\{n_jx_j\}$ 有界，故 $\{n_jTx_j\}$ 有界。拓扑有界性保证有界集乘以趋零标量时一致趋于零，因此 $Tx_j=n_j^{-1}(n_jTx_j)\to0$。度量空间上线性映射的序列连续等价于连续，故结论成立。
\end{proof}
\begin{corollary}[补充：可度量局部凸空间是囿空间]
每个可度量的局部凸空间都是囿空间。
\end{corollary}
\begin{proof}
对在有界集上有界的半范数 $p$，将它看作线性映射 $X\to X/\ker p$ 的目标范数。此映射保有界集，上一定理（目标赋范）保证连续，故 $p$ 连续，再引用定理 \ref{thm:y1-bornological}。
\end{proof}
囿空间关注“有界性足以检测连续性”；它与以后出现的桶空间不同。后者的定义要求所有闭的吸收绝对凸集都是邻域。两者的假设和用途需要分别检查。

\section{广义函数和广义导数}\label{sec:yosida-01-08}
\noindent\textit{原书位置：中译本第 48--54 页；英文版第 46--52 页。}

\subsection{测试函数空间的有界集与囿性}
\begin{theorem}[测试空间有界集的完整刻画]\label{thm:y1-D-bounded}
$B\subset\Test(\Omega)$ 有界，当且仅当存在共同紧集 $K\Subset\Omega$，使 $B\subset\Test_K$，且对每个 $m$，
\[
 \sup_{\varphi\in B}p_{K,m}(\varphi)<\infty.
\]
$\Test(\Omega)$ 是囿空间。
\end{theorem}
\begin{proof}
若无共同紧集，与定理 \ref{thm:y1-D-sequence} 同样选出 $\varphi_j\in B$ 及无域内聚点的非零取值点 $x_j$。半范数
\[
 q(\psi)=\sum_{j=1}^\infty j\frac{|\psi(x_j)|}{|\varphi_j(x_j)|}
\]
各层限制为有限和，故连续；却有 $q(\varphi_j)\ge j$，与有界性矛盾。各阶导数的一致有界由全局连续半范数 $\sup_\Omega|\partial^\alpha\psi|$ 得到。反向给定任意连续半范数 $q$，其在 $\Test_K$ 上由有限个 $p_{K,m}$ 控制，所以在 $B$ 上有界，命题 \ref{prop:y1-bounded} 即给出充分性。

最后，若半范数 $q$ 在 $\Test$ 的有界集上有界，其在每个 $\Test_K$ 上也保有界性，因为包含映射连续。$\Test_K$ 可度量且局部凸，故囿，$q|_{\Test_K}$ 连续。归纳极限判别给出 $q$ 连续，于是 $\Test$ 囿。这个证明无需把不可度量的 $\Test$ 当作度量空间。
\end{proof}

\subsection{定义、局部有限阶估计与典型分布}
\begin{definition}[广义函数或分布]
$\Omega$ 上的分布是连续线性泛函 $T:\Test(\Omega)\to\mathbb K$；全体分布记作 $\Dist(\Omega)$。配对记为 $T(\varphi)=\Pair{T}{\varphi}$。
\end{definition}
\begin{theorem}[分布的局部有限阶判别]\label{thm:y1-distribution-estimate}
线性泛函 $T$ 是分布，当且仅当对每个 $K\Subset\Omega$，存在 $m\in\N_0$ 与 $C_K>0$，使
\[
 |\Pair{T}{\varphi}|\le C_Kp_{K,m}(\varphi)\quad(\varphi\in\Test_K).
\]
它也等价于 $T$ 在 $\Test$ 的每个有界集上有界。
\end{theorem}
\begin{proof}
连续性限制到 $\Test_K$，用第 \ref{sec:yosida-01-06} 节的半范数估计；$p_{K,m}$ 随 $m$ 增大，所以有限最大值可用一个较高阶半范数替代。反向由归纳极限判别得到。最后一条由 $\Test$ 的囿性以及半范数 $|T(\cdot)|$ 的判别得到。
\end{proof}
阶数 $m$ 可随紧集变化；一般分布未必有一个对整个 $\Omega$ 通用的有限阶数。

\begin{example}[局部可积函数、测度和 Dirac 分布]
对 $f\in L^1_{\Loc}(\Omega)$，令 $T_f(\varphi)=\int_\Omega f\varphi\,dx$。其估计为 $|T_f(\varphi)|\le\|f\|_{L^1(K)}p_{K,0}(\varphi)$。若复 Radon 测度 $\nu$ 的全变差在紧集上有限，则 $T_\nu(\varphi)=\int\varphi\,d\nu$ 同样定义分布。特别地，$\delta_a(\varphi)=\varphi(a)$ 是集中于 $a$ 的 Dirac 分布。仅有 $\sigma$-有限性不能代替局部有限全变差条件。
\end{example}

\begin{proposition}[局部可积函数到分布的嵌入是单射]\label{prop:y1-regular-injective}
若 $T_f=0$，则 $f=0$ 几乎处处。因而 $T_f=T_g$ 等价于 $f=g$ 几乎处处。
\end{proposition}
\begin{proof}
给定任意 $V\Subset\Omega$ 的开集，在包含 $\overline V$ 的紧邻域内取足够小的 $\varepsilon$，则对 $x\in V$，$y\mapsto\rho_\varepsilon(x-y)$ 是测试函数，所以 $f*\rho_\varepsilon(x)=0$。局部可积函数的正则化在 $L^1(V)$ 中趋于 $f$。这一逼近可由局部截断化为 $L^1(\R^n)$，再用
\[
 \|f*\rho_\varepsilon-f\|_1
 \le\int\rho_\varepsilon(h)\|f(\cdot-h)-f\|_1\,dh
\]
及 $L^1$ 平移连续性证明；后者由简单函数的矩形逼近得到。因此 $f=0$ 在 $V$ 中几乎处处，取可数穷竭完成证明。这避免将局部可积的 $f$ 在整个无界域上当成有限测度。
\end{proof}

\subsection{广义导数与跳跃项}
\begin{definition}[分布导数]
对任意多重指标 $\alpha$，定义
\[
 \Pair{\partial^\alpha T}{\varphi}
 =(-1)^{|\alpha|}\Pair{T}{\partial^\alpha\varphi}.
\]
\end{definition}
这是分布，因为测试函数的微分保持紧支集，且提高有限阶估计中的阶数。各分布都能任意次求导，偏导相互交换。若 $f\in C^1(\Omega)$，分部积分给出 $\partial_jT_f=T_{\partial_jf}$；测试函数在边界附近为零，所以不产生边界项。

\begin{example}[Heaviside 函数与分段光滑函数]
对 $H=1_{(0,\infty)}$，
\[
 \Pair{H'}{\varphi}=-\int_0^\infty\varphi'(x)\,dx=\varphi(0),\qquad H'=\delta_0.
\]
若 $f$ 在有限个点 $a_j$ 外分段 $C^1$，具有有限单侧极限，且分段导数局部可积，则
\[
 \partial T_f=T_{f'_{\mathrm{pw}}}
 +\sum_j\bigl(f(a_j+)-f(a_j-)\bigr)\delta_{a_j}.
\]
证明是在各光滑区间分部积分，相邻端点项合并为右值减左值。点值本身任意更改不影响公式。特别地，$(|x|)''=2\delta_0$。
\end{example}

\begin{proposition}[边界与 Green 公式]\label{prop:y1-boundary-distribution}
设 $U\Subset\R^n$ 有光滑边界，$f\in C^2(\overline U)$，并把 $f$ 延零到 $U$ 外。取外法向 $n$，则
\[
 \partial_jT_{f1_U}=T_{(\partial_jf)1_U}
 -f n_j\,dS\big|_{\partial U},
\]
且对测试函数 $\varphi$，
\[
 \Pair{\Delta T_{f1_U}}{\varphi}
 =\int_U\Delta f\,\varphi\,dx
 +\int_{\partial U}\bigl(f\partial_n\varphi-(\partial_nf)\varphi\bigr)\,dS.
\]
\end{proposition}
\begin{proof}
散度定理给出 $-\int_U f\partial_j\varphi=\int_U\partial_jf\,\varphi-\int_{\partial U}fn_j\varphi$。Green 第二恒等式给出 $\int_U f\Delta\varphi=\int_U\Delta f\,\varphi+\int_{\partial U}(f\partial_n\varphi-\partial_nf\,\varphi)$。原书取内法向时两条边界符号相应反转；必须先固定法向约定。
\end{proof}

\subsection{光滑乘子与广义 Leibniz 公式}
对 $a\in C^\infty(\Omega)$ 定义 $\Pair{aT}{\varphi}=\Pair{T}{a\varphi}$。Leibniz 公式给出 $p_{K,m}(a\varphi)\le C_{a,K,m}p_{K,m}(\varphi)$，所以定义的是分布，且
\[
 \partial_j(aT)=(\partial_ja)T+a\partial_jT.
\]
两边作用于 $\varphi$ 后，只需将 $\partial_j(a\varphi)$ 展开即得。

\begin{theorem}[Hörmander 的广义 Leibniz 公式]\label{thm:y1-general-leibniz}
为与后面的 Fourier 记号衔接，此处临时记 $D_j=(1/i)\partial_j$。若 $P(\xi)$ 是多项式，$a$ 光滑，$T$ 是分布，则
\[
 P(D)(aT)=\sum_\beta\frac{D^\beta a}{\beta!}
 (\partial_\xi^\beta P)(D)T.
\]
右边是有限和。
\end{theorem}
\begin{proof}
先对单项式 $P(\xi)=\xi^\alpha$ 逐次应用乘积求导，得
\[
 D^\alpha(aT)=\sum_{\beta\le\alpha}\binom\alpha\beta
 (D^\beta a)D^{\alpha-\beta}T.
\]
而 $\partial_\xi^\beta\xi^\alpha/\beta!=\binom\alpha\beta\xi^{\alpha-\beta}$，所以单项式情形成立。再以线性组合得到一般多项式。这补足了原书利用指数函数和多项式 Taylor 恒等式辨认系数的论证。
\end{proof}
一般两个分布之间没有自然乘法：例如 $\delta_0^2$ 不由以上定义产生。对光滑乘子的要求是保证乘积仍把每个测试函数送入测试空间。

\begin{exercise}[补充：Dirac 分布与乘子]
证明在 $\R$ 上，$x\delta_0=0$、$x\delta_0'=-\delta_0$，并求 $a\delta_0'$ 的一般表达式。
\end{exercise}
\begin{solution}
分别作用于 $\varphi$。$x\delta_0$ 给出零；$x\delta_0'$ 给出 $-(x\varphi)'(0)=-\varphi(0)$；一般有 $a\delta_0'=a(0)\delta_0'-a'(0)\delta_0$。由此可见不能将 Dirac 分布误当成只在一点有非零数值的普通函数。
\end{solution}
进一步的拓扑、限制、运算与例子参见第 \ref{ch:distributions}、\ref{ch:operations} 章。

\section{B-空间和 F-空间}\label{sec:yosida-01-09}
\noindent\textit{原书位置：中译本第 54--58 页；英文版第 52--56 页。}

\subsection{完备性与术语}
\begin{definition}[B-空间、F-空间与 Hilbert 空间]
在平移不变距离下，每个 Cauchy 序列均收敛的空间称为完备。完备赋范空间称为 Banach 空间，原书简称 B-空间；完备拟赋范空间称为 F-空间。完备准 Hilbert 空间称为 Hilbert 空间。现代通常把“Fréchet 空间”限定为完备、可度量的局部凸空间；一般 F-空间不必局部凸。
\end{definition}
收敛序列总是 Cauchy；极限唯一，因为 $d(x,y)\le d(x,x_j)+d(x_j,y)$。完备性要求极限落在空间内，而不只落在某个更大的函数集合内。

\begin{theorem}[光滑函数空间与固定支集测试空间完备]\label{thm:y1-smooth-complete}
$\mathcal E^k(\Omega)$（$k$ 有限或无穷）以及 $\Test_K(\Omega)$ 都是 Fréchet 空间。
\end{theorem}
\begin{proof}
对 $\mathcal E^k$ 中的 Cauchy 序列 $f_j$，每个允许的 $\partial^\alpha f_j$ 在每个紧集上一致 Cauchy，故有连续局部极限 $g_\alpha$，不同紧集上的极限相容。需要证明 $g_\alpha$ 真是 $g_0$ 的导数。在域内任意小闭长方体上，对 $|\alpha|<k$ 用微积分基本定理写
\[
 \partial^\alpha f_j(x+te_l)-\partial^\alpha f_j(x)
 =\int_0^t\partial^{\alpha+e_l}f_j(x+se_l)\,ds.
\]
一致极限允许交换积分与极限，得同样的 $g$ 公式；再由连续性求导，得到 $\partial_lg_\alpha=g_{\alpha+e_l}$。归纳即有 $g_0\in C^k$，并在所有半范数下收敛。无限阶情形对每个有限阶重复即可。若原序列均支集于 $K$，域外 $\Omega\setminus K$ 上所有函数及其极限为零，故极限仍属于 $\Test_K$。这补足了“导数各有极限，所以极限光滑”中的关键相容性步骤。
\end{proof}

\subsection{可积函数空间的完备性}
\begin{theorem}[Riesz--Fischer 完备性]\label{thm:y1-Lp-complete}
对任意测度空间，$L^p$（$1\le p\le\infty$）都是 Banach 空间。若 $1\le p<\infty$，每个 $L^p$-Cauchy 序列都有一个几乎处处收敛于其范数极限的子列。
\end{theorem}
\begin{proof}
先设 $p<\infty$。从 Cauchy 序列抽出 $f_{j_k}$，使
\[
 \|f_{j_{k+1}}-f_{j_k}\|_p\le2^{-k}.
\]
令 $g_N=|f_{j_1}|+\sum_{k=1}^N|f_{j_{k+1}}-f_{j_k}|$。Minkowski 保证 $\|g_N\|_p\le\|f_{j_1}\|_p+1$；单调收敛或 Fatou 引理给出极限 $g\in L^p$。因此 $g<\infty$ 几乎处处，差分级数绝对收敛，定义几乎处处有限的
\[
 f=f_{j_1}+\sum_{k=1}^\infty(f_{j_{k+1}}-f_{j_k}).
\]
$|f|\le g$ 保证 $f\in L^p$。再次用 Fatou 与 Minkowski，
\[
 \|f-f_{j_k}\|_p\le\sum_{l\ge k}2^{-l}\longrightarrow0.
\]
原序列由 Cauchy 性和这一收敛子列也趋于 $f$。

对于 $p=\infty$，同样选择子列使差分范数不超过 $2^{-k}$。去掉可数个零测集后，每个差分逐点不超过该上界，故差分级数在余集上一致收敛。极限本质有界，且 $L^\infty$ 距离的尾项不超过 $\sum_{l\ge k}2^{-l}$；再由 Cauchy 性处理原序列。甚至原序列在去掉一个共同零测集后可取一致收敛的代表元。
\end{proof}
$C_b(S)$ 也完备：一致 Cauchy 序列逐点取极限，一致收敛保留连续性和有界性。由此 $c_0,c$ 是 Banach 空间；$\ell^p$ 是 $L^p$ 的计数测度特例。因而 $L^2,\ell^2$ 是 Hilbert 空间，定理 \ref{thm:y1-hardy} 又使 $H^2$ 成为 Hilbert 空间。

\begin{proposition}[补充：有界算子空间的完备性]
若 $Y$ 是 Banach 空间，则 $\mathcal L(X,Y)$ 在算子范数下是 Banach 空间，不要求 $X$ 完备。
\end{proposition}
\begin{proof}
算子范数 Cauchy 序列 $T_j$ 对每个 $x$ 使 $T_jx$ 在 $Y$ 中 Cauchy，定义 $Tx=\lim T_jx$。逐点极限保留线性；$\|T_j\|$ 统一有界，所以 $T$ 有界。给定 $\varepsilon$，$j,l\ge N$ 时 $\|(T_j-T_l)x\|\le\varepsilon\|x\|$，令 $l\to\infty$，再在单位球取上确界，得到 $\|T_j-T\|\le\varepsilon$。
\end{proof}

\subsection{Bergman 空间、测度收敛空间与全体数列}
\begin{proposition}[Bergman 空间完备]
$A^2(G)$ 是 $L^2(G)$ 的闭子空间，因而是 Hilbert 空间。
\end{proposition}
\begin{proof}
全纯函数在 $\overline B(z,r)\subset G$ 上的圆周平均或 Taylor 展开给出
\[
 |f(z)|^2\le\frac1{\pi r^2}\int_{B(z,r)}|f(w)|^2\,dA(w).
\]
例如展开 $f(z+\rho e^{i\theta})=\sum a_j\rho^je^{ij\theta}$，角积分消掉异次项，再对 $\rho\,d\rho$ 积分，常数项就给出右边至少 $|a_0|^2$。在任意紧集 $K\Subset G$ 上可取统一 $r>0$，所以 $L^2$-Cauchy 序列的全纯代表元在 $K$ 上一致 Cauchy。其局部一致极限 $f$ 全纯，可由 Cauchy 积分公式取极限证明。与此同时 $L^2$ 完备性给出范数极限 $g$，两种极限在每个紧子域上几乎处处一致，故 $g$ 有全纯代表 $f$。所以闭性成立。
\end{proof}

\begin{proposition}[测度收敛空间完备]
若 $\mu(S)<\infty$，则第 \ref{sec:yosida-01-04} 节的 $L^0(S,\mu)$ 是 F-空间。
\end{proposition}
\begin{proof}
从距离 Cauchy 序列抽出子列 $f_{j_k}$，使
\[
 \mu\{|f_{j_{k+1}}-f_{j_k}|>2^{-k}\}<2^{-k}.
\]
尾部异常集之并的测度不超过 $\sum_{l\ge k}2^{-l}$，趋于零；因此几乎处处只有有限次异常，差分级数的尾部绝对值受 $\sum2^{-l}$ 控制，得到有限极限 $f$。有限测度上的几乎处处收敛推出测度收敛，故子列在距离中趋于 $f$；原序列由 Cauchy 性得到同样结论。
\end{proof}
原书的空间 $(s)$，即所有实或复数列 $\mathbb K^{\N}$，在
\[
 q(x)=\sum_{j=1}^\infty2^{-j}\frac{|x_j|}{1+|x_j|}
\]
下是 F-空间且局部凸：Cauchy 性等价于各坐标 Cauchy，逐坐标取极限后，有限首项与小尾项论证给出距离收敛。其拓扑就是逐坐标收敛的乘积拓扑。
\begin{pitfall}[补充核查：原书的另一计数公式]
原书还列出 $\arctan\inf_{\varepsilon>0}(\varepsilon+\#\{j:|x_j|>\varepsilon\})$ 作为替代公式。若它按未加权的普通计数理解，则不能在\emph{全体}数列上作为拟范数：取 $x_j=j$，任意有限阈值以上的坐标数都是无穷，所以任意 $a\ne0$ 时该表达式在 $ax$ 上恒为 $\pi/2$，不满足 $a\to0$ 的数乘连续性。本稿因此采用上面的加权求和公式来处理全体数列；这里没有把有缺陷的替代式作为定理使用。
\end{pitfall}

\subsection{Sobolev 空间及其完备性}
\begin{definition}[弱导数与 Sobolev 空间]
对整数 $k\ge0$ 和 $1\le p<\infty$，定义
\[
 W^{k,p}(\Omega)=\{f\in L^p(\Omega):\partial^\alpha T_f=T_{g_\alpha},\
 g_\alpha\in L^p(\Omega),\ |\alpha|\le k\}.
\]
由分布嵌入单射，$g_\alpha$ 唯一，记为弱导数 $\partial^\alpha f$。范数为
\[
 \|f\|_{k,p}=\left(\sum_{|\alpha|\le k}\|\partial^\alpha f\|_p^p\right)^{1/p}.
\]
当 $p=2$，记 $H^k(\Omega)=W^{k,2}(\Omega)$，取第 \ref{sec:yosida-01-05} 节的导数内积。
\end{definition}
\begin{theorem}[Sobolev 空间完备]\label{thm:y1-sobolev-complete}
$W^{k,p}(\Omega)$ 是 Banach 空间；$H^k(\Omega)$ 是 Hilbert 空间。
\end{theorem}
\begin{proof}
若 $f_j$ 在 Sobolev 范数下 Cauchy，则各 $\partial^\alpha f_j$ 在 $L^p$ 中 Cauchy，设极限为 $g_\alpha$，特别地 $f_j\to g_0=f$。对任意 $\varphi\in\Test$，Hölder（$p=1$ 时用上确界）允许在
\[
 \int\partial^\alpha f_j\,\varphi
 =(-1)^{|\alpha|}\int f_j\partial^\alpha\varphi
\]
两边取极限，得到 $\int g_\alpha\varphi=(-1)^{|\alpha|}\int f\partial^\alpha\varphi$。所以 $g_\alpha$ 正是 $f$ 的弱导数，$f\in W^{k,p}$，全部有限项范数同时收敛。这是“函数及其导数独立取极限”后必须验证的相容性。$p=2$ 时范数由内积生成，结合完备性即得 Hilbert 性。
\end{proof}

\begin{exercise}[补充：不完备的自然函数空间]
令 $c_{00}$ 为有限支集数列空间，取 $\ell^2$ 范数。构造一个不在 $c_{00}$ 中收敛的 Cauchy 序列，并确定其在 $\ell^2$ 中的极限。
\end{exercise}
\begin{solution}
取 $x^{(N)}=(1,1/2,\ldots,1/N,0,\ldots)$。差的平方范数为相应的 $\sum j^{-2}$ 尾部，故 Cauchy。极限 $(1/j)_{j\ge1}$ 属于 $\ell^2$，但不属于 $c_{00}$；若在 $c_{00}$ 内有极限，坐标连续性会迫使它等于该无限支集数列，矛盾。
\end{solution}

\section{完备化}\label{sec:yosida-01-10}
\noindent\textit{原书位置：中译本第 58--61 页；英文版第 56--59 页。}

\subsection{由 Cauchy 序列构造完备化}
\begin{theorem}[拟赋范空间的完备化]\label{thm:y1-completion}
每个拟赋范空间 $X$ 都能线性等距嵌入某个 F-空间 $\widehat X$，且像稠密；该完备化在固定原空间嵌入的线性等距同构意义下唯一。若原拟范数是范数，则 $\widehat X$ 是 Banach 空间；若原范数来自内积，则其完备化是 Hilbert 空间。
\end{theorem}
\begin{proof}
令 $\mathcal C$ 为 $X$ 的全部 Cauchy 序列，定义 $x_j\sim y_j$ 当且仅当 $q(x_j-y_j)\to0$。三角不等式证明它是等价关系。定义
\[
 [x_j]+[y_j]=[x_j+y_j],\qquad a[x_j]=[ax_j],\qquad
 \widehat q([x_j])=\lim_jq(x_j).
\]
反三角不等式保证实数序列 $q(x_j)$ Cauchy，故极限存在。加法与拟范数不依赖代表元；固定数乘的连续性保证 $[ax_j]$ 仍为 Cauchy 序列且不依赖代表元。正定性、对称性及三角不等式直接继承。

\emph{补充检查数乘条件。} 一个 Cauchy 序列 $x_j$ 满足 $q(ax_j)\to0$ 随 $a\to0$ 对 $j$ 一致。确实，数乘在 $(0,0)$ 联合连续，所以给定误差可先选小的向量邻域和标量邻域，使两者乘积小；再把 $x_j$ 的尾部写成某个固定 $x_N$ 加上该小向量，有限个首项另行控制。因此 $\widehat q(a[x_j])\to0$。对固定 $a$，源空间的数乘连续性给出：$q(x)<\delta$ 时 $q(ax)<\varepsilon/2$。若 $\widehat q(\xi)<\delta/2$，其代表元最终满足 $q(x_j)<\delta$，故 $\widehat q(a\xi)\le\varepsilon/2<\varepsilon$。于是 $\widehat q$ 满足两条数乘条件。这避免未经论证交换两重极限。

嵌入 $Jx=[x,x,\ldots]$ 线性等距。对 $\xi=[x_j]$，有 $\widehat q(\xi-Jx_N)=\lim_jq(x_j-x_N)\to0$，所以像稠密。若 $\xi_k$ 在 $\widehat X$ 中 Cauchy，由稠密性取 $z_k\in X$ 使 $\widehat q(Jz_k-\xi_k)<2^{-k}$。三角不等式使 $z_k$ 在 $X$ 中 Cauchy，令 $\xi=[z_k]$；由刚证明的稠密逼近及 Cauchy 性，$Jz_k\to\xi$，从而 $\xi_k\to\xi$，得完备性。

若另有稠密等距嵌入 $J':X\to Z$，$Z$ 完备，则映射 $[x_j]\mapsto\lim_jJ'x_j$ 良定义、线性等距且满射：任何 $z\in Z$ 可由原空间的像逼近。它固定 $X$，所以唯一。范数齐次性以及平行四边形恒等式均可在极限中继承。
\end{proof}

\begin{proposition}[补充：算子的延拓性质]\label{prop:y1-extension}
设 $Y$ 是 Banach 空间，$T:X\to Y$ 为赋范空间上的有界线性算子，则存在唯一有界线性延拓 $\widehat T:\widehat X\to Y$，且 $\|\widehat T\|=\|T\|$。
\end{proposition}
\begin{proof}
定义 $\widehat T([x_j])=\lim_jTx_j$。估计 $\|Tx_j-Tx_l\|\le\|T\|\|x_j-x_l\|$ 保证极限存在及代表元无关，并给出 $\|\widehat T\|\le\|T\|$。限制于原空间得到反向不等式。稠密性保证唯一性。
\end{proof}

\subsection{Sobolev 空间的完备化解释}
令 $H_0^k(\Omega)$ 为 $C_c^\infty(\Omega)$ 在 $H^k$ 范数下的完备化。第 \ref{sec:yosida-01-09} 节的完备性与弱导数相容性证明说明，它可等距识别为
\[
 H_0^k(\Omega)=\overline{C_c^\infty(\Omega)}^{\,H^k(\Omega)}.
\]
因此它是 $H^k(\Omega)$ 的闭子空间。原书还从 $\widehat H^k$ 出发作抽象完备化；本稿将现代弱导数定义的空间直接记为 $H^k$，以免“经典导数空间”和“已完备的 Sobolev 空间”混淆。对任意开集，$H_0^k$ 的定义不需要边界光滑；“零边界值”的迹解释则需要另加相应条件。

\begin{lemma}[补充：$L^p$ 平移连续与正则化逼近]\label{lem:y1-Lp-translation}
若 $1\le p<\infty$，$f\in L^p(\R^n)$，则 $\|f(\cdot-h)-f\|_p\to0$ 随 $h\to0$。因而 $f*\rho_\varepsilon\to f$ 在 $L^p$ 中。
\end{lemma}
\begin{proof}
先对有界矩形的示性函数证明：平移后对称差的体积趋零，$L^p$ 距离的 $p$ 次方正是该体积。有限个这样的简单函数亦成立。一般 $L^p$ 函数可先截断大小和远方尾部，再用简单函数及有限测度集的矩形逼近，在 $L^p$ 中逼近为上述函数。平移等距，于是
\[
 \|f(\cdot-h)-f\|_p
 \le2\|f-g\|_p+\|g(\cdot-h)-g\|_p
\]
完成一般情形。卷积逼近由 Jensen 不等式和 Fubini 得
\[
 \|f*\rho_\varepsilon-f\|_p^p
 \le\int\rho_\varepsilon(h)\|f(\cdot-h)-f\|_p^p\,dh.
\]
核支集收缩至零，右边趋零。
\end{proof}

\begin{theorem}[全空间上的光滑紧支集稠密性]\label{thm:y1-sobolev-density}
对每个整数 $k\ge0$ 与 $1\le p<\infty$，$C_c^\infty(\R^n)$ 在 $W^{k,p}(\R^n)$ 中稠密。特别地
\[
 H_0^k(\R^n)=H^k(\R^n)=W^{k,2}(\R^n).
\]
\end{theorem}
\begin{proof}
取 $\chi\in C_c^\infty$ 在单位球上为 $1$，令 $\chi_R(x)=\chi(x/R)$。Leibniz 公式给出
\[
 \partial^\alpha(\chi_Rf)-\partial^\alpha f
 =(\chi_R-1)\partial^\alpha f
 +\sum_{0<\beta\le\alpha}\binom\alpha\beta
 (\partial^\beta\chi_R)\partial^{\alpha-\beta}f.
\]
第一项的 $L^p$ 范数由 $\partial^\alpha f$ 的可积尾部趋零；其余项的系数不超过 $C_\beta R^{-|\beta|}$，且支集位于远方环带，故也趋零。因此 $\chi_Rf\to f$ 在 Sobolev 范数下，先把问题化为紧支集。

固定 $R$ 后令 $f_{R,\varepsilon}=(\chi_Rf)*\rho_\varepsilon$。它光滑且紧支集。利用分布导数的定义，将微分移入卷积核，再移到 $\chi_Rf$ 的弱导数，得到
\[
 \partial^\alpha f_{R,\varepsilon}
 =(\partial^\alpha(\chi_Rf))*\rho_\varepsilon.
\]
引理 \ref{lem:y1-Lp-translation} 保证每个有限阶导数在 $L^p$ 下逼近；先选大 $R$，再选小 $\varepsilon$，使总误差任意小。
\end{proof}

\begin{example}[补充：在有边界的域中两空间可以不同]
在 $(0,1)$ 上，常数 $1\in H^1(0,1)$，却不属于 $H_0^1(0,1)$。若 $\varphi_j\in C_c^\infty(0,1)$ 在 $H^1$ 中趋于 $1$，令 $v_j=\varphi_j-1$。对任意 $v\in C^1([0,1])$，由 $v(0)=v(x)-\int_0^xv'(t)\,dt$ 并对 $x$ 积分，得
\[
 |v(0)|\le\|v\|_{L^2(0,1)}+\|v'\|_{L^2(0,1)}.
\]
应用于 $v_j$，右侧趋零，左侧恒为 $1$，矛盾。这也说明将全空间的密度结论直接搬到任意有边界域会出错。
\end{example}

\section{B-空间的商空间}\label{sec:yosida-01-11}
\noindent\textit{原书位置：中译本第 61--62 页；英文版第 59--60 页。}

\begin{definition}[商范数]
设 $X$ 赋范，$M\subset X$ 为闭线性子空间。对商空间中的陪集 $x+M$ 定义
\[
 \|x+M\|_{X/M}=\inf_{m\in M}\|x+m\|=\operatorname{dist}(x,M).
\]
规范映射记为 $Q:X\to X/M$，$Qx=x+M$。
\end{definition}
\begin{proposition}[商范数、商拓扑与开映射]\label{prop:y1-quotient-norm}
该定义不依赖代表元，是商空间上的范数；其拓扑正是商拓扑。$Q$ 线性连续，$\|Q\|\le1$，且 $Q$ 是开映射。
\end{proposition}
\begin{proof}
代表元变化只是把 $m$ 平移，故下确界不变。齐次性显然，三角性由对两个下确界选取任意精度的近似代表再相加得到。范数为零时 $x$ 在 $M$ 的闭包中，闭性给出 $x\in M$，故正定。

有 $\|Qx\|\le\|x\|$。更准确地，
\[
 Q(B_X(x,r))=B_{X/M}(Qx,r).
\]
一向由压缩估计；反向若 $\|\xi-Qx\|<r$，下确界的定义允许选取该陪集中的代表 $z$，使 $\|z\|<r$，于是 $\xi=Q(x+z)$。这证明 $Q$ 开，也证明范数拓扑等于商拓扑。
\end{proof}
若 $M$ 不闭，同一公式仍是半范数，其零空间为 $\overline M/M$；所以闭性不是一个仅为方便而附加的条件。

\begin{theorem}[Banach 空间的闭子空间商仍完备]\label{thm:y1-quotient-complete}
若 $X$ 为 Banach 空间且 $M$ 闭，则 $X/M$ 为 Banach 空间。
\end{theorem}
\begin{proof}
从商空间的 Cauchy 序列 $\xi_j$ 抽出子列，使 $\|\xi_{j_{k+1}}-\xi_{j_k}\|<2^{-k-1}$。对每个差陪集选代表 $u_k\in X$，使
\[
 Qu_k=\xi_{j_{k+1}}-\xi_{j_k},\qquad\|u_k\|<2^{-k}.
\]
选 $x_1$ 满足 $Qx_1=\xi_{j_1}$。级数 $x=x_1+\sum_ku_k$ 在 Banach 空间中绝对收敛。部分和 $s_k=x_1+\sum_{l<k}u_l$ 满足 $Qs_k=\xi_{j_k}$，因此 $\xi_{j_k}\to Qx$。原序列由 Cauchy 性也趋于 $Qx$。这里仅需要近似达到下确界的代表，完全不需要最近点存在。
\end{proof}

图 \ref{fig:y1-quotient} 给出商范数的有限维模型；证明中的代表元只须近似达到下确界。
\begin{figure}[htbp]
\centering
\begin{tikzpicture}[fa diagram]
  \draw[fa axis] (-.7,0)--(5.6,0) node[right] {$M$};
  \draw[fa main] (-.5,2)--(5,2) node[right] {$x+M$};
  \node[fa dot,label=above:$x$] at (3.7,2) {};
  \node[fa dot,label=below:$0$] at (1,0) {};
  \draw[fa accent] (1,0)--(1,2) node[midway,left] {$d$};
  \draw[fa main,dashed] (1,0)--(3.7,2);
  \draw[fa main,dashed] (1,0)--(4.7,2);
  \node[fa note] at (3.2,-1) {$\|Qx\|=\inf_{m\in M}\|x+m\|$};
  \node[fa box] at (8.5,1) {商空间把整条陪集\\视为同一个点 $Qx$};
\end{tikzpicture}
\caption{商范数在 $\R^2$ 欧氏范数中的模型：在同一陪集中改变代表元，再取距原点的距离下确界。本例中 $d$ 由垂足达到；一般 Banach 空间的闭子空间未必有最近点，正文完备性证明只选择近似代表。}
\label{fig:y1-quotient}
\end{figure}



\begin{proposition}[补充：经商空间分解算子]
若有界线性算子 $T:X\to Y$ 满足 $M\subset\ker T$，则存在唯一有界线性算子 $\widetilde T:X/M\to Y$，满足 $T=\widetilde TQ$，并且 $\|\widetilde T\|=\|T\|$。
\end{proposition}
\begin{proof}
定义 $\widetilde T(x+M)=Tx$。$M\subset\ker T$ 保证代表元无关；由 $\|Tx\|=\|T(x+m)\|\le\|T\|\|x+m\|$ 对 $m$ 取下确界，得到 $\|\widetilde T\|\le\|T\|$。反向来自 $T=\widetilde TQ$ 与 $\|Q\|\le1$。满射性保证唯一性。
\end{proof}

\begin{exercise}[补充：商空间的具体计算]
在 $C([0,1])$ 上取上确界范数，令 $M=\{f:f(0)=0\}$。证明 $C([0,1])/M$ 与标量域等距同构。
\end{exercise}
\begin{solution}
评价映射 $f\mapsto f(0)$ 的核为 $M$，故诱导双射。所有代表的范数至少为 $|f(0)|$；而陪集中有常函数 $f(0)$，其范数恰为 $|f(0)|$，故商范数就是 $|f(0)|$。
\end{solution}

\section{单位分解}\label{sec:yosida-01-12}
\noindent\textit{原书位置：中译本第 63--64 页；英文版第 60--62 页。}

\subsection{局部有限覆盖的作用}
一族集合称为局部有限，若每个点都有邻域只与其中有限多个集合相交。在开集 $\Omega$ 中，对闭支集组成的局部有限族，每个紧集只与有限多个支集相交；因此相应的函数和在局部是有限和，可以逐项求导。这是把局部定义拼成全局对象的基本机制。

\begin{theorem}[从属任意开覆盖的光滑单位分解]\label{thm:y1-partition}
设 $\Omega\subset\R^n$ 开，$\{G_i\}_{i\in I}$ 是其开覆盖。存在可数函数族 $\{\theta_j\}_{j\ge1}\subset C_c^\infty(\Omega)$，满足
\[
 0\le\theta_j\le1,\qquad \supp\theta_j\subset G_{i(j)},\qquad
 \sum_j\theta_j(x)=1\quad(x\in\Omega),
\]
且支集族在 $\Omega$ 中局部有限。这里允许多个 $j$ 从属于同一个覆盖元素。
\end{theorem}
\begin{proof}
取紧集穷竭 $K_m\subset\operatorname{int}K_{m+1}$，置 $K_0=K_{-1}=\varnothing$。紧环层
\[
 A_m=K_m\setminus\operatorname{int}K_{m-1}
\]
包含于开环层 $W_m=\operatorname{int}K_{m+1}\setminus K_{m-2}$。对每个 $x\in A_m$，选择覆盖元素 $G_i$ 及小球 $B(x,2r)$，使闭球包含于 $W_m\cap G_i$。有限个小球 $B(x,r)$ 覆盖 $A_m$。在每个大球中取非负光滑紧支集函数 $b_j$，在对应小球上严格为正。

所有环层的这些函数组成可数族，且支集在 $\Omega$ 中局部有限：固定紧集包含于某个 $K_N$，而 $m\ge N+2$ 的支集在 $K_N$ 外，剩下层各只有有限个函数。函数
\[
 b(x)=\sum_jb_j(x)
\]
因局部有限而光滑，因小球覆盖各环层而处处为正。令 $\theta_j=b_j/b$，则正性、求和为一和支集条件全部成立。每个 $b_j$ 的支集紧包含于 $\Omega$，所以 $\theta_j$ 也可延零成 $\R^n$ 上的光滑紧支集函数；分母可能在域边界变小不影响这一点，因为每个单独支集与边界保持正距离。
\end{proof}
这给出了原书“散开覆盖”的标准局部有限含义，同时确保局部有限的是\emph{支集}而非仅非零开集，消除了归一化与求导中的潜在缺口。

\begin{corollary}[光滑截断函数]\label{cor:y1-cutoff}
若 $K\Subset\Omega$、$O\subset\Omega$ 开且含 $K$，则存在 $\chi\in\Test(\Omega)$，$0\le\chi\le1$，在 $K$ 的某个邻域上为 $1$，且 $\supp\chi\subset O$。
\end{corollary}
\begin{proof}
用有限个闭球包含于 $O$ 的小球覆盖 $K$，取在小球上为正的紧支集光滑函数，相加得 $b$，其支集紧包含于 $O$，且 $\min_K b=c>0$。取光滑实函数 $h$，在 $t\le c/2$ 上为零，在 $t\ge3c/4$ 上为一，并介于零与一之间。$\chi=h\circ b$ 即满足要求。
\end{proof}
特别地，可以取 $\chi_m$ 在 $K_m$ 附近为一，于是 $\chi_mf\to f$ 在 $\mathcal E(\Omega)$ 中成立，说明 $\Test$ 在 $\mathcal E$ 中稠密。这个稠密性不是说 $\Test$ 与 $\mathcal E$ 的拓扑相同。

\subsection{补充：分布的限制与粘合}
若 $U\subset\Omega$ 开，$\Test(U)$ 的函数延零到 $\Omega$ 后仍是测试函数，因此可定义 $T|_U$。下面的局部原则使后面支集的定义成立。
\begin{proposition}[分布的局部唯一性与粘合]\label{prop:y1-gluing}
若分布 $T$ 在开集族 $U_i$ 上均为零，则它在 $\bigcup_iU_i$ 上为零。更一般地，若 $T_i\in\Dist(U_i)$ 在交集上相容，则存在唯一 $T\in\Dist(\bigcup_iU_i)$，其各限制等于 $T_i$。
\end{proposition}
\begin{proof}
先取从属于 $U_i$ 的单位分解 $\theta_j$。对支集于并集的 $\varphi$，只有有限个 $\theta_j\varphi$ 非零，因此 $\varphi=\sum_j\theta_j\varphi$。若各局部限制为零，对有限和应用 $T$ 即得零。

一般情形定义
\[
 T(\varphi)=\sum_jT_{i(j)}(\theta_j\varphi).
\]
对固定 $\Test_K$ 只有有限项；每项由光滑乘子连续性及局部有限阶估计控制，所以 $T$ 是分布。若 $\supp\varphi\subset U_i$，各 $\theta_j\varphi$ 的紧支集在 $U_i\cap U_{i(j)}$ 中，相容性使和等于 $\sum_jT_i(\theta_j\varphi)=T_i(\varphi)$。局部零结论保证唯一性，因此也保证结果与单位分解选择无关。
\end{proof}
专题部分第 \ref{ch:distributions} 章给出限制与粘合的进一步应用。

\section{具有紧支集的广义函数}\label{sec:yosida-01-13}
\noindent\textit{原书位置：中译本第 64--67 页；英文版第 62--65 页。}

\subsection{支集及光滑函数对偶}
\begin{definition}[分布支集]
称 $T$ 在开集 $U$ 上为零，若它对所有支集于 $U$ 的测试函数取零。由命题 \ref{prop:y1-gluing}，所有这种开集的并仍是零区域。其补集称为 $\supp T$，即 $T$ 的支集。
\end{definition}
支集是相对闭集。若 $\varphi$ 的紧支集与 $\supp T$ 不相交，则 $T(\varphi)=0$。由定义可直接得
\[
 \supp(\partial^\alpha T)\subset\supp T,\qquad
 \supp(aT)\subset\supp a\cap\supp T.
\]
例如 $\supp\delta_a=\{a\}$，而非零光滑函数生成的分布支集就是其函数支集；后一事实利用局部可积嵌入单射。

$\mathcal E(\Omega)=C^\infty(\Omega)$ 的有界集由每个紧集上的每阶导数一致有界刻画，直接来自半范数有界判别。它是可度量局部凸空间，故囿，因此其线性泛函连续等价于在所有有界集上有界。

\begin{theorem}[紧支集分布等同于光滑函数的连续对偶]\label{thm:y1-Eprime}
每个紧支集分布 $T$ 唯一延拓为 $\mathcal E(\Omega)$ 上的连续线性泛函。反过来，每个 $L\in\mathcal E(\Omega)'$ 限制于 $\Test$ 后都是紧支集分布。这给出自然的线性对应
\[
 \CompactDist(\Omega)=\{T\in\Dist(\Omega):\supp T\Subset\Omega\}.
\]
连续性还等价于存在一个 $K\Subset\Omega$、整数 $m\ge0$ 和常数 $C$，使
\[
 |L(f)|\le C\max_{|\alpha|\le m}\sup_K|\partial^\alpha f|
 \quad(f\in C^\infty(\Omega)).
\]
\end{theorem}
\begin{proof}
设 $\supp T=F$ 紧，取 $\chi\in\Test$ 在 $F$ 附近为一，定义 $L(f)=T(\chi f)$。若换用另一个 $\chi_1$，$(\chi-\chi_1)f$ 是紧支集测试函数且在 $F$ 附近为零，故 $T$ 对它取零，定义与截断选择无关。由 $T$ 在固定紧集 $K=\supp\chi$ 上的有限阶估计及 Leibniz 公式，得到上述对 $f$ 的估计，因此 $L$ 连续。它限制到 $\Test$ 等于 $T$；$\Test$ 在 $\mathcal E$ 中稠密，故连续延拓唯一。

反向，连续性给出有限个 $p_{K_j,m_j}$ 控制 $|L(f)|$，取有限个紧集的并及最大的阶数得所列单一估计。于是 $L|_{\Test}$ 是分布，并在 $\Omega\setminus K$ 上为零，故支集包含于 $K$。估计的充分性显然。注意这里“一种延拓”指的是延拓\emph{泛函}唯一；实现它的截断函数可以有很多个。
\end{proof}
每个紧支集分布因此具有全局有限阶估计，但估计中的紧集可以稍大于其真实支集。Dirac 分布及其导数可作用于任何光滑函数；一般非紧支集分布则不能未经条件直接作用于常函数 $1$。

\subsection{点支集分布的完整结构}
\begin{theorem}[点支集分布是 Dirac 导数的有限和]\label{thm:y1-point-support}
若 $\supp T\subset\{a\}$，则存在 $m<\infty$ 及唯一系数 $c_\alpha$，使
\[
 T=\sum_{|\alpha|\le m}c_\alpha\partial^\alpha\delta_a.
\]
只在需要时选取足够大的 $m$；补零后各系数唯一。
\end{theorem}
\begin{proof}
平移后设 $a=0$。由定理 \ref{thm:y1-Eprime}，存在固定紧邻域 $K$ 上的 $m$ 阶估计。若 $\partial^\alpha f(0)=0$ 对 $|\alpha|\le m$ 成立，取 $\chi\in\Test$ 在零附近为一，令 $\chi_\varepsilon(x)=\chi(x/\varepsilon)$，其支集在充分小时落于 $K$。支集性质给出 $T(f)=T(\chi_\varepsilon f)$。

Taylor 余项保证对 $|\gamma|\le m$，在 $|x|\le C\varepsilon$ 上，
\[
 |\partial^\gamma f(x)|\le C_f\varepsilon^{m+1-|\gamma|}.
\]
Leibniz 公式中的一项为 $\varepsilon^{-|\beta|}(\partial^\beta\chi)(x/\varepsilon)\partial^{\alpha-\beta}f(x)$，因此对 $|\alpha|\le m$，
\[
 \sup_K|\partial^\alpha(\chi_\varepsilon f)|
 \le C_f'\varepsilon^{m+1-|\alpha|}\le C_f''\varepsilon
 \quad(0<\varepsilon\le1).
\]
有限阶估计使 $T(f)=0$。这展开了截断函数求导产生负幂后，Taylor 余项仍能抵消它们的理由。

一般 $f$ 减去 $m$ 阶 Taylor 多项式 $P_mf=\sum_{|\alpha|\le m}\partial^\alpha f(0)x^\alpha/\alpha!$ 后满足上述零喷射条件，所以 $T(f)=T(P_mf)$。这里 $T$ 已延拓到 $\mathcal E$，可以合法作用于多项式。由于
\[
 (\partial^\alpha\delta_0)(f)=(-1)^{|\alpha|}\partial^\alpha f(0),
\]
取 $c_\alpha=(-1)^{|\alpha|}T(x^\alpha/\alpha!)$ 即得表示。最后以在零附近等于单项式的紧支集测试函数分别检验各喷射，得到系数唯一。
\end{proof}
有限多个点支集的分布也是各点 Dirac 导数的有限和：取互不相交的小邻域和相应截断，将分布分解成各点支集分布后应用定理。点支集结论不能把“有限阶”换成任意无穷导数级数，连续性正是截断为有限喷射的原因。

\begin{exercise}[补充：点支集分布的阶数]
证明在 $\R$ 上 $\delta_0^{(m)}$ 的阶数恰为 $m$：它受 $m$ 阶测试半范数控制，但 $m\ge1$ 时不能由 $m-1$ 阶半范数控制。
\end{exercise}
\begin{solution}
上界由 $|\delta_0^{(m)}(\varphi)|=|\varphi^{(m)}(0)|$ 给出。选 $\eta\in\Test$ 使 $\eta^{(m)}(0)\ne0$，令 $\varphi_\varepsilon(x)=\varepsilon^{m-1}\eta(x/\varepsilon)$。支集保持在一个固定紧邻域中，所有至 $m-1$ 阶导数统一有界，而 $|\varphi_\varepsilon^{(m)}(0)|=\varepsilon^{-1}|\eta^{(m)}(0)|\to\infty$。因此不存在低阶估计。
\end{solution}

\section{广义函数的直接积}\label{sec:yosida-01-14}
\noindent\textit{原书位置：中译本第 67--70 页；英文版第 65--67 页。}

\subsection{分离变量函数的稠密性}
记 $u\otimes v$ 为函数 $(x,y)\mapsto u(x)v(y)$。有限个这样的函数的线性组合构成测试空间的代数张量积；它不是集合的笛卡尔积，也不是把同一变量中的两个分布相乘。

\begin{theorem}[分离变量的测试函数逼近]\label{thm:y1-tensor-density}
给定 $\varphi\in\Test(\R^{n+m})$，存在
\[
 \varphi_j(x,y)=\sum_{l=1}^{N_j}u_{jl}(x)v_{jl}(y),\qquad
 u_{jl}\in\Test(\R^n),\ v_{jl}\in\Test(\R^m),
\]
使 $\varphi_j\to\varphi$ 在 $\Test(\R^{n+m})$ 中成立。即有限分离变量函数的和在完整测试拓扑下稠密。
\end{theorem}
\begin{proof}
展开原书的高斯正则化证明。令 $d=n+m$、$z=(x,y)$，定义
\[
 \Phi_t(z)=(4\pi t)^{-d/2}\int_{\R^d}
 \varphi(w)\exp\left(-\frac{|z-w|^2}{4t}\right)\,dw,
 \qquad t>0.
\]
对每个导数 $\partial^\alpha$，分部积分将微分从核移到 $\varphi$，得到 $\partial^\alpha\Phi_t=G_t*\partial^\alpha\varphi$。换元 $w=z+2\sqrt t\,s$ 后，积分为归一化高斯下 $\partial^\alpha\varphi(z+2\sqrt t\,s)$ 的平均。分成 $|s|\le R$ 与 $|s|>R$：前一部分因紧支集光滑函数的全域一致连续性而随 $t\downarrow0$ 一致逼近，后一部分由可积高斯尾部先取大 $R$ 控制。因此所有有限阶导数一致趋于 $\partial^\alpha\varphi$。

积分中的 $|z-w|^2$ 在复变量延拓时应理解为 $\sum(z_l-w_l)^2$，不用复共轭。$w$ 的积分区间固定紧集，所以核对复 $z$ 是整函数并在复紧集上有统一可积控制；故 $\Phi_t$ 是 $d$ 变量整函数。其多元 Taylor 部分和为多项式，在每个复紧集上一致收敛；由 Cauchy 积分公式，各有限阶导数也一致收敛。选择 $t_j\downarrow0$ 及足够高阶多项式 $P_j$，使在 $|z|\le j$ 上、$|\alpha|\le j$ 时，$\partial^\alpha(P_j-\Phi_{t_j})$ 的绝对值小于 $1/j$。于是 $P_j\to\varphi$ 在 $\mathcal E(\R^d)$ 中。

选固定 $a\in\Test(\R^n)$、$b\in\Test(\R^m)$，分别在 $\supp\varphi$ 的两坐标投影附近为一。令 $\varphi_j=a(x)b(y)P_j(x,y)$。它们的支集全部在固定紧集 $\supp a\times\supp b$ 中，且 Leibniz 公式给出所有导数一致趋于 $a(x)b(y)\varphi=\varphi$，所以是 $\Test$ 收敛。每个多项式是有限个 $x^\alpha y^\beta$ 的和，乘上 $a,b$ 后即是所需分离变量和。固定截断而非随 $j$ 逃向无穷的截断，是最后一步的要点。
\end{proof}
同一结论适用于 $\Test(\Omega\times\Lambda)$：先用有限个紧包含于乘积域的矩形邻域及单位分解分解一个给定测试函数，再对每片选择支集留在对应坐标域中的截断。因支集紧，只需有限片，每片的上述逼近可以合并。

\subsection{带参数作用、直积及 Fubini 性质}
\begin{lemma}[分布作用后的参数函数仍是测试函数]\label{lem:y1-parameter}
设 $S\in\Dist(\Lambda)$，$\varphi\in\Test(\Omega\times\Lambda)$。则
\[
 \psi(x)=\Pair{S}{\varphi(x,\cdot)}
\]
属于 $\Test(\Omega)$，且 $\partial_x^\alpha\psi(x)=\Pair{S}{\partial_x^\alpha\varphi(x,\cdot)}$。若 $\supp\varphi\subset K\times L$ 且 $S$ 在 $\Test_L$ 上受 $s$ 阶估计控制，则
\[
 \max_{|\alpha|\le r}\sup_K|\partial_x^\alpha\psi|
 \le C_S\max_{\substack{|\alpha|\le r\\|\beta|\le s}}
 \sup_{K\times L}|\partial_x^\alpha\partial_y^\beta\varphi|.
\]
\end{lemma}
\begin{proof}
$x$ 在坐标投影支集外时，$\varphi(x,\cdot)=0$，所以 $\psi$ 有紧支集。对一阶差商，Taylor 公式给出
\[
 \frac{\varphi(x+he_j,y)-\varphi(x,y)}h-\partial_{x_j}\varphi(x,y)
 \longrightarrow0
\]
在任意有限阶的 $y$ 导数半范数下一致成立，且 $y$ 支集固定于 $L$。$S$ 连续允许取此极限，得到导数公式；逐次重复证明所有阶。对导数公式应用 $S$ 的有限阶估计即得最后的上界。它还证明带参数作用把测试空间的有界集送入有界集。
\end{proof}

\begin{theorem}[分布直积的存在、唯一性与交换作用顺序]\label{thm:y1-tensor}
对 $T\in\Dist(\Omega)$、$S\in\Dist(\Lambda)$，存在唯一 $T\otimes S\in\Dist(\Omega\times\Lambda)$，使
\[
 \Pair{T\otimes S}{u(x)v(y)}=\Pair{T}{u}\Pair{S}{v}.
\]
对一般测试函数，具有分布版 Fubini 公式
\[
 \Pair{T\otimes S}{\varphi}
 =\Pair{T_x}{\Pair{S_y}{\varphi(x,y)}}
 =\Pair{S_y}{\Pair{T_x}{\varphi(x,y)}}.
\]
\end{theorem}
\begin{proof}
由上一引理，先作用 $S$ 再作用 $T$ 有意义，并且在固定 $K\times L$ 上有
\[
 |T_x(S_y\varphi)|\le C_TC_S
 \max_{\substack{|\alpha|\le r\\|\beta|\le s}}
 \sup_{K\times L}|\partial_x^\alpha\partial_y^\beta\varphi|.
\]
每个一般紧集都包含在两个坐标投影紧集的乘积内，故局部有限阶判别证明它是分布。反向作用顺序同样定义分布。两者对分离变量函数都取值 $T(u)S(v)$，故对其有限线性组合相同。定理 \ref{thm:y1-tensor-density} 的稠密性及两个分布的连续性证明它们在全部测试函数上相同，并同时证明满足所列乘积条件的分布唯一。
\end{proof}
交换公式是“交换作用于不同变量的顺序”；若把 $T\otimes S$ 与 $S\otimes T$ 看作不同次序乘积域上的分布，还须配合坐标交换映射。原书以 $\times$ 记直接积，此处用 $\otimes$ 突出独立变量与张量含义。

\begin{proposition}[补充：支集、微分与普通函数的对应]
分布直积满足
\[
 \supp(T\otimes S)=\supp T\times\supp S,
\]
以及
\[
 \partial_x^\alpha\partial_y^\beta(T\otimes S)
 =(\partial^\alpha T)\otimes(\partial^\beta S).
\]
若 $f\in L^1_{\Loc}(\Omega)$、$g\in L^1_{\Loc}(\Lambda)$，则 $T_f\otimes T_g=T_{f(x)g(y)}$。
\end{proposition}
\begin{proof}
若一点在 $\supp T\times\supp S$ 外，可取一个矩形邻域，其中某个因子分布为零，直积在此邻域为零，得支集包含关系。反向若 $a\in\supp T$、$b\in\supp S$，任意包含 $(a,b)$ 的邻域都有较小矩形 $U\times V$。由支集定义存在支集于 $U,V$ 的 $u,v$，使 $T(u)\ne0$、$S(v)\ne0$，于是直积在 $uv$ 上非零，故 $(a,b)$ 属于支集。微分公式逐次使用导数定义和 Fubini 公式。普通函数情形中，固定乘积紧集上的绝对可积性由 $\int_K|f|\int_L|g|<\infty$ 保证，经典 Fubini 给出对应公式。
\end{proof}
例如 $\delta_a\otimes\delta_b=\delta_{(a,b)}$，且 $\delta_a'\otimes\delta_b$ 作用于 $\varphi$ 给出 $-\partial_x\varphi(a,b)$。这都合法，仍没有定义同一变量中的 $\delta_a^2$。更深入的卷积与张量应用见第 \ref{ch:convolution} 章。

\begin{exercise}[补充：直积与微分的检验]
在 $\R^2$ 上令 $W=H(x)\otimes\delta_0(y)$。写出 $W(\varphi)$，计算 $\partial_xW$ 与 $\partial_yW$，并确定其支集。
\end{exercise}
\begin{solution}
$W(\varphi)=\int_0^\infty\varphi(x,0)\,dx$。因此 $\partial_xW=\delta_0(x)\otimes\delta_0(y)=\delta_{(0,0)}$，$\partial_yW=H(x)\otimes\delta_0'(y)$，后者作用于测试函数为 $-\int_0^\infty\partial_y\varphi(x,0)\,dx$。支集为 $[0,\infty)\times\{0\}$。
\end{solution}


\section*{本章参考文献与阅读衔接}
\phantomsection
\addcontentsline{toc}{section}{本章参考文献与阅读衔接}
本章十四节的基本材料来源于 \cite{yosida-zh,yosida-en} 第一章。原书章末还指引读者阅读 Bourbaki、Grothendieck、Köthe、Banach、Dunford--Schwartz 及 Hille--Phillips 关于局部凸空间与 Banach 空间的著作，以及 Schwartz、Gelfand--Shilov、Hörmander 与 Friedman 关于分布的著作。本稿实际使用的分布参考书见 \cite{schwartz,gelfand}；上述其余著作仅保留为原书的阅读线索，未作为已核查的直接引文。

第 \ref{ch:yosida-02} 章将从完备性推出一致有界、开映射与闭图像定理；第 \ref{ch:yosida-03} 章继续发展 Hilbert 空间的投影与表示。广义函数的进一步例子、卷积、Fourier 变换及偏微分方程应用，可接续阅读第 \ref{ch:test}--\ref{ch:weighted} 章。本章中分布的基本定义与证明仍独立保留，以保证主线阅读不依赖专题部分。
